\RequirePackage{fix-cm}
%DIF LATEXDIFF DIFFERENCE FILE
%DIF DEL itt-emotion_second_half_V1.tex                     Wed Feb 10 22:51:19 2021
%DIF ADD itt-emotion_second_half_V3_service_corrected.tex   Tue Feb 16 02:42:59 2021
%
%\documentclass{svjour3}                     % onecolumn (standard format)
%\documentclass[smallcondensed]{svjour3}     % onecolumn (ditto)
\documentclass[smallextended,natbib]{svjour3}       % onecolumn (second format)
% Fold 1
  % Fold 2
    % Fold 3
      % Fold 4
        %\documentclass[twocolumn]{svjour3}          % twocolumn
        %
        \smartqed  % flush right qed marks, e.g. at end of proof

        % For the ORCID icon
        \usepackage{scalerel}
        \usepackage{tikz}
        \usetikzlibrary{svg.path}

        \definecolor{orcidlogocol}{HTML}{A6CE39}
        \tikzset{
          orcidlogo/.pic={
            \fill[orcidlogocol] svg{M256,128c0,70.7-57.3,128-128,128C57.3,256,0,198.7,0,128C0,57.3,57.3,0,128,0C198.7,0,256,57.3,256,128z};
            \fill[white] svg{M86.3,186.2H70.9V79.1h15.4v48.4V186.2z}
                         svg{M108.9,79.1h41.6c39.6,0,57,28.3,57,53.6c0,27.5-21.5,53.6-56.8,53.6h-41.8V79.1z M124.3,172.4h24.5c34.9,0,42.9-26.5,42.9-39.7c0-21.5-13.7-39.7-43.7-39.7h-23.7V172.4z}
                         svg{M88.7,56.8c0,5.5-4.5,10.1-10.1,10.1c-5.6,0-10.1-4.6-10.1-10.1c0-5.6,4.5-10.1,10.1-10.1C84.2,46.7,88.7,51.3,88.7,56.8z};
          }
        }

        \newcommand\orcidicon[1]{\href{https://orcid.org/#1}{\mbox{\scalerel*{
        \begin{tikzpicture}[yscale=-1,transform shape]
        \pic{orcidlogo};
        \end{tikzpicture}
        }{|}}}}

        \usepackage{lineno, hyperref}
        \usepackage{array, pbox}
        \usepackage{mathtools}
        \usepackage{caption}
        \usepackage{subcaption}
        \usepackage{graphicx}

        % To use chinese characters
        \usepackage{CJKutf8}

        % To add an appendix page
        \usepackage[page, title]{appendix}

        % To insert commas into large numbers automatically
        \usepackage{siunitx}

        % To add research questions
        \usepackage{ntheorem}
        \theoremseparator{:}
        \newtheorem{rsq}{Research Question}

        % To add subrsq theorems
        \makeatletter
        \newcounter{subrsq} 
        \let\savedc@rsq\c@rsq
        \newenvironment{subrsq}
         {%
          \setcounter{subrsq}{0}%
          \stepcounter{rsq}%
          \edef\saved@rsq{\thersq}% Save the current value of rsq
          \let\c@rsq\c@subrsq     % Now rsq is subrsq
          \renewcommand{\thersq}{\saved@rsq\alph{rsq}}%
         }
         {}
        \newcommand{\normrsq}{%
          \let\c@rsq\savedc@rsq % revert to the old one
          \renewcommand\thersq{\arabic{rsq}}%
        } 
        \makeatother


        % To add the envelope symbol in Corresponding Author using % \Letter
        \usepackage{marvosym}

        % to add superscripts to mark the institute of authors
        % patch \maketitle:
        \newcommand*{\affaddr}[1]{#1} % No op here. Customize it for different styles.
        \newcommand*{\affmark}[1][*]{\textsuperscript{#1}}

        % For table formatting
        \usepackage{multirow}
        \usepackage[normalem]{ulem}
        \useunder{\uline}{\ul}{}

        % For landscape tables
        \usepackage{lscape}

        \usepackage{bibentry}

        \modulolinenumbers[5]

        \journalname{Information Technology \& Tourism}

        \bibliographystyle{spbasic}
%DIF PREAMBLE EXTENSION ADDED BY LATEXDIFF
%DIF UNDERLINE PREAMBLE %DIF PREAMBLE
\RequirePackage[normalem]{ulem} %DIF PREAMBLE
\RequirePackage{color}\definecolor{RED}{rgb}{1,0,0}\definecolor{BLUE}{rgb}{0,0,1} %DIF PREAMBLE
\providecommand{\DIFadd}[1]{{\protect\color{blue}\uwave{#1}}} %DIF PREAMBLE
\providecommand{\DIFdel}[1]{{\protect\color{red}\sout{#1}}}                      %DIF PREAMBLE
%DIF SAFE PREAMBLE %DIF PREAMBLE
\providecommand{\DIFaddbegin}{} %DIF PREAMBLE
\providecommand{\DIFaddend}{} %DIF PREAMBLE
\providecommand{\DIFdelbegin}{} %DIF PREAMBLE
\providecommand{\DIFdelend}{} %DIF PREAMBLE
\providecommand{\DIFmodbegin}{} %DIF PREAMBLE
\providecommand{\DIFmodend}{} %DIF PREAMBLE
%DIF FLOATSAFE PREAMBLE %DIF PREAMBLE
\providecommand{\DIFaddFL}[1]{\DIFadd{#1}} %DIF PREAMBLE
\providecommand{\DIFdelFL}[1]{\DIFdel{#1}} %DIF PREAMBLE
\providecommand{\DIFaddbeginFL}{} %DIF PREAMBLE
\providecommand{\DIFaddendFL}{} %DIF PREAMBLE
\providecommand{\DIFdelbeginFL}{} %DIF PREAMBLE
\providecommand{\DIFdelendFL}{} %DIF PREAMBLE
\newcommand{\DIFscaledelfig}{0.5}
%DIF HIGHLIGHTGRAPHICS PREAMBLE %DIF PREAMBLE
\RequirePackage{settobox} %DIF PREAMBLE
\RequirePackage{letltxmacro} %DIF PREAMBLE
\newsavebox{\DIFdelgraphicsbox} %DIF PREAMBLE
\newlength{\DIFdelgraphicswidth} %DIF PREAMBLE
\newlength{\DIFdelgraphicsheight} %DIF PREAMBLE
% store original definition of \includegraphics %DIF PREAMBLE
\LetLtxMacro{\DIFOincludegraphics}{\includegraphics} %DIF PREAMBLE
\newcommand{\DIFaddincludegraphics}[2][]{{\color{blue}\fbox{\DIFOincludegraphics[#1]{#2}}}} %DIF PREAMBLE
\newcommand{\DIFdelincludegraphics}[2][]{% %DIF PREAMBLE
\sbox{\DIFdelgraphicsbox}{\DIFOincludegraphics[#1]{#2}}% %DIF PREAMBLE
\settoboxwidth{\DIFdelgraphicswidth}{\DIFdelgraphicsbox} %DIF PREAMBLE
\settoboxtotalheight{\DIFdelgraphicsheight}{\DIFdelgraphicsbox} %DIF PREAMBLE
\scalebox{\DIFscaledelfig}{% %DIF PREAMBLE
\parbox[b]{\DIFdelgraphicswidth}{\usebox{\DIFdelgraphicsbox}\\[-\baselineskip] \rule{\DIFdelgraphicswidth}{0em}}\llap{\resizebox{\DIFdelgraphicswidth}{\DIFdelgraphicsheight}{% %DIF PREAMBLE
\setlength{\unitlength}{\DIFdelgraphicswidth}% %DIF PREAMBLE
\begin{picture}(1,1)% %DIF PREAMBLE
\thicklines\linethickness{2pt} %DIF PREAMBLE
{\color[rgb]{1,0,0}\put(0,0){\framebox(1,1){}}}% %DIF PREAMBLE
{\color[rgb]{1,0,0}\put(0,0){\line( 1,1){1}}}% %DIF PREAMBLE
{\color[rgb]{1,0,0}\put(0,1){\line(1,-1){1}}}% %DIF PREAMBLE
\end{picture}% %DIF PREAMBLE
}\hspace*{3pt}}} %DIF PREAMBLE
} %DIF PREAMBLE
\LetLtxMacro{\DIFOaddbegin}{\DIFaddbegin} %DIF PREAMBLE
\LetLtxMacro{\DIFOaddend}{\DIFaddend} %DIF PREAMBLE
\LetLtxMacro{\DIFOdelbegin}{\DIFdelbegin} %DIF PREAMBLE
\LetLtxMacro{\DIFOdelend}{\DIFdelend} %DIF PREAMBLE
\DeclareRobustCommand{\DIFaddbegin}{\DIFOaddbegin \let\includegraphics\DIFaddincludegraphics} %DIF PREAMBLE
\DeclareRobustCommand{\DIFaddend}{\DIFOaddend \let\includegraphics\DIFOincludegraphics} %DIF PREAMBLE
\DeclareRobustCommand{\DIFdelbegin}{\DIFOdelbegin \let\includegraphics\DIFdelincludegraphics} %DIF PREAMBLE
\DeclareRobustCommand{\DIFdelend}{\DIFOaddend \let\includegraphics\DIFOincludegraphics} %DIF PREAMBLE
\LetLtxMacro{\DIFOaddbeginFL}{\DIFaddbeginFL} %DIF PREAMBLE
\LetLtxMacro{\DIFOaddendFL}{\DIFaddendFL} %DIF PREAMBLE
\LetLtxMacro{\DIFOdelbeginFL}{\DIFdelbeginFL} %DIF PREAMBLE
\LetLtxMacro{\DIFOdelendFL}{\DIFdelendFL} %DIF PREAMBLE
\DeclareRobustCommand{\DIFaddbeginFL}{\DIFOaddbeginFL \let\includegraphics\DIFaddincludegraphics} %DIF PREAMBLE
\DeclareRobustCommand{\DIFaddendFL}{\DIFOaddendFL \let\includegraphics\DIFOincludegraphics} %DIF PREAMBLE
\DeclareRobustCommand{\DIFdelbeginFL}{\DIFOdelbeginFL \let\includegraphics\DIFdelincludegraphics} %DIF PREAMBLE
\DeclareRobustCommand{\DIFdelendFL}{\DIFOaddendFL \let\includegraphics\DIFOincludegraphics} %DIF PREAMBLE
%DIF LISTINGS PREAMBLE %DIF PREAMBLE
\RequirePackage{listings} %DIF PREAMBLE
\RequirePackage{color} %DIF PREAMBLE
\lstdefinelanguage{DIFcode}{ %DIF PREAMBLE
%DIF DIFCODE_UNDERLINE %DIF PREAMBLE
  moredelim=[il][\color{red}\sout]{\%DIF\ <\ }, %DIF PREAMBLE
  moredelim=[il][\color{blue}\uwave]{\%DIF\ >\ } %DIF PREAMBLE
} %DIF PREAMBLE
\lstdefinestyle{DIFverbatimstyle}{ %DIF PREAMBLE
	language=DIFcode, %DIF PREAMBLE
	basicstyle=\ttfamily, %DIF PREAMBLE
	columns=fullflexible, %DIF PREAMBLE
	keepspaces=true %DIF PREAMBLE
} %DIF PREAMBLE
\lstnewenvironment{DIFverbatim}{\lstset{style=DIFverbatimstyle}}{} %DIF PREAMBLE
\lstnewenvironment{DIFverbatim*}{\lstset{style=DIFverbatimstyle,showspaces=true}}{} %DIF PREAMBLE
%DIF END PREAMBLE EXTENSION ADDED BY LATEXDIFF

\begin{document}

\title{Differences in Chinese and Western tourists faced with Japanese hospitality: A natural language processing approach}
% Fold 1
  % Fold 2
    % Fold 3
      % Fold 4
        \titlerunning{Differences in Chinese and Western tourists faced with Japanese hospitality: ...}

        \author{REDACTED}

        \authorrunning{E. Alemán Carreón et al.}

        \institute{REDACTED}

        \date{Received: date / Accepted: date}
        % The correct dates will be entered by the editor

        \maketitle

\section{Results}\label{results}

  \subsection{\DIFdelbegin \DIFdel{Experiment }\DIFdelend \DIFaddbegin \DIFadd{Experimental }\DIFaddend results and \DIFdelbegin \DIFdel{answering }\DIFdelend \DIFaddbegin \DIFadd{answers to }\DIFaddend research questions}

    Our research questions were \DIFdelbegin \DIFdel{about two things. In }\DIFdelend \DIFaddbegin \DIFadd{related to two issues. Based on }\DIFaddend research questions \ref{rsq:hospitality} and \ref{rsq:hospitality_both}, \DIFdelbegin \DIFdel{we decide that }\DIFdelend the objective of this study \DIFaddbegin \DIFadd{was }\DIFaddend to determine how Chinese and Western tourists \DIFdelbegin \DIFdel{interact with }\DIFdelend \DIFaddbegin \DIFadd{respond to }\DIFaddend the \textit{omotenashi} culture\DIFdelbegin \DIFdel{influenced }\DIFdelend \DIFaddbegin \DIFadd{, which influences }\DIFaddend hospitality and service in Japan, and \DIFdelbegin \DIFdel{how are they different in their perceptions in this matter}\DIFdelend \DIFaddbegin \DIFadd{the differences in their perception of this hospitality}\DIFaddend . We observed the top-ranking positive factors for Chinese tourists across different price ranges\DIFaddbegin \DIFadd{, as shown }\DIFaddend in Table \ref{tab:freq_res_pos}, and specifically the word ``\begin{CJK}{UTF8}{gbsn}不错\end{CJK} (not bad)'' and its pairings\DIFaddbegin \DIFadd{, as shown }\DIFaddend in Table \ref{tab:adj_zh_pos}. \DIFdelbegin \DIFdel{From these observations , we can infer }\DIFdelend \DIFaddbegin \DIFadd{These observations revealed }\DIFaddend that, while service, cleanliness, and breakfast \DIFdelbegin \DIFdel{are }\DIFdelend \DIFaddbegin \DIFadd{were }\DIFaddend praised in most hotels, location \DIFdelbegin \DIFdel{is usually placed above it }\DIFdelend \DIFaddbegin \DIFadd{usually placed higher }\DIFaddend in importance on the pairings. \DIFdelbegin \DIFdel{When we see }\DIFdelend \DIFaddbegin \DIFadd{From }\DIFaddend the rest of the factors \DIFaddbegin \DIFadd{that are }\DIFaddend lower on the list, we \DIFdelbegin \DIFdel{see }\DIFdelend \DIFaddbegin \DIFadd{inferred }\DIFaddend that the list \DIFdelbegin \DIFdel{is }\DIFdelend \DIFaddbegin \DIFadd{was }\DIFaddend more populated with hard attributes \DIFdelbegin \DIFdel{like }\DIFdelend \DIFaddbegin \DIFadd{such as }\DIFaddend location and transportation availability across different price ranges. \DIFdelbegin \DIFdel{From the }\DIFdelend \DIFaddbegin \DIFadd{The }\DIFaddend negative keyword usages in Table \ref{tab:freq_res_neg} \DIFdelbegin \DIFdel{, there are }\DIFdelend \DIFaddbegin \DIFadd{indicate that there were }\DIFaddend complaints about the lack of a \DIFdelbegin \DIFdel{Chinese friendly environment. However, most complaints are also }\DIFdelend \DIFaddbegin \DIFadd{Chinese-friendly environment, but there were more complaints }\DIFaddend about hard attributes such as the building's age and the distance from other convenient spots. \DIFdelbegin \DIFdel{Nevertheless, the most complained about aspect is }\DIFdelend \DIFaddbegin \DIFadd{However, most complaints were about }\DIFaddend the price of the hotel\DIFdelbegin \DIFdel{. Surprisingly, }\DIFdelend \DIFaddbegin \DIFadd{, which included }\DIFaddend all of the price ranges\DIFdelbegin \DIFdel{have this negative keyword at the top of the list, suggesting that it is the main concern to }\DIFdelend \DIFaddbegin \DIFadd{; therefore, the price was the main concern for }\DIFaddend Chinese customers with different travel purposes.

    \DIFdelbegin \DIFdel{On the other hand, the }\DIFdelend \DIFaddbegin \DIFadd{The }\DIFaddend word ``staff'' is the second or third in \DIFdelbegin \DIFdel{the }\DIFdelend \DIFaddbegin \DIFadd{all the price ranges in the }\DIFaddend lists of satisfaction factors in \DIFdelbegin \DIFdel{English-written reviews in all the price ranges}\DIFdelend \DIFaddbegin \DIFadd{reviews written in English }\DIFaddend . This word is followed by a few other keywords \DIFaddbegin \DIFadd{that are }\DIFaddend lower in the top 10 list, such as ``helpful'' or ``friendly\DIFdelbegin \DIFdel{''.When we look at }\DIFdelend \DIFaddbegin \DIFadd{.'' From }\DIFaddend the pairings of the \DIFdelbegin \DIFdel{top-ranked }\DIFdelend \DIFaddbegin \DIFadd{top-ranking }\DIFaddend keyword ``good'' in Table \ref{tab:adj_en_pos}, we \DIFdelbegin \DIFdel{find that customers mostly }\DIFdelend \DIFaddbegin \DIFadd{realize that customers mainly }\DIFaddend praise the location, service, breakfast, or English availability. \DIFdelbegin \DIFdel{When we look at }\DIFdelend \DIFaddbegin \DIFadd{From }\DIFaddend the negative keyword ``poor'' and its pairings in Table \ref{tab:adj_en_neg}, we \DIFdelbegin \DIFdel{see that  it is also service-related concepts that }\DIFdelend \DIFaddbegin \DIFadd{deduce that  }\DIFaddend the Western tourists are \DIFdelbegin \DIFdel{disappointed with. With these results, we can observe that }\DIFdelend \DIFaddbegin \DIFadd{also disappointed with service-related concepts. Therefore, }\DIFaddend both Chinese and English-speaking tourists in Japan have different priorities. However, both populations consider the hotel's location and \DIFdelbegin \DIFdel{transport availability }\DIFdelend \DIFaddbegin \DIFadd{transportation availability nearby }\DIFaddend (subways and trains) \DIFdelbegin \DIFdel{nearby as secondary but still essential points in their satisfactionwith a hotel}\DIFdelend \DIFaddbegin \DIFadd{as secondary yet essential points to their satisfaction}\DIFaddend . The Chinese customers are primarily satisfied with the room quality in \DIFaddbegin \DIFadd{terms of the }\DIFaddend spaciousness and cleanliness \DIFdelbegin \DIFdel{and }\DIFdelend \DIFaddbegin \DIFadd{as well as }\DIFaddend the service of breakfast.

    In contrast, the English-speaking customers \DIFdelbegin \DIFdel{are }\DIFdelend \DIFaddbegin \DIFadd{were }\DIFaddend easily upset by \DIFdelbegin \DIFdel{any lack of cleanliness }\DIFdelend \DIFaddbegin \DIFadd{uncleanliness }\DIFaddend and smoke smell from cigarettes. \DIFdelbegin \DIFdel{Surprisingly, the cigarette smell is }\DIFdelend \DIFaddbegin \DIFadd{Cigarette smell was }\DIFaddend an issue even in the \DIFdelbegin \DIFdel{middle to }\DIFdelend \DIFaddbegin \DIFadd{middle- and }\DIFaddend high-class hotels\DIFdelbegin \DIFdel{above }\DIFdelend \DIFaddbegin \DIFadd{, of which the rooms were priced at more than }\DIFaddend 30,000 yen per night. \DIFdelbegin \DIFdel{However, }\DIFdelend \DIFaddbegin \DIFadd{For hotels with rooms priced }\DIFaddend above 50,000 yen per night, \DIFdelbegin \DIFdel{this problem seems }\DIFdelend \DIFaddbegin \DIFadd{however, this problem seemed }\DIFaddend to disappear from the list of top 10 concerns. \DIFdelbegin \DIFdel{Old and dated buildings seem to be }\DIFdelend \DIFaddbegin \DIFadd{Furthermore, the age of the buildings was }\DIFaddend a concern for both populations. \DIFdelbegin \DIFdel{On the positive side}\DIFdelend \DIFaddbegin \DIFadd{Nevertheless}\DIFaddend , for all price ranges considered, English-speaking tourists \DIFdelbegin \DIFdel{value }\DIFdelend \DIFaddbegin \DIFadd{valued }\DIFaddend staff friendliness over room quality when considering their satisfaction. In contrast, Chinese tourists \DIFdelbegin \DIFdel{consider }\DIFdelend \DIFaddbegin \DIFadd{considered }\DIFaddend location and transportation more often.

    We \DIFdelbegin \DIFdel{also can }\DIFdelend \DIFaddbegin \DIFadd{can also }\DIFaddend observe some keywords that are not considered by their counterparts. For example, English-speaking customers mentioned tobacco smell in many reviews. However, it was not statistically identified as a problem for their Chinese counterparts. On the other hand, \DIFdelbegin \DIFdel{while it appears }\DIFdelend \DIFaddbegin \DIFadd{although they appear }\DIFaddend in both English and Chinese lists, references to ``\begin{CJK}{UTF8}{gbsn}购物\end{CJK} (shopping)'' are more common in the Chinese lists across hotels of 15,000 yen to 200,000 yen per night. Meanwhile, the term ``shopping'' \DIFdelbegin \DIFdel{only appears in the }\DIFdelend \DIFaddbegin \DIFadd{appeared solely in the top 10 positive keywords list for English speakers who stayed in rooms priced }\DIFaddend 20,000\DIFdelbegin \DIFdel{to }\DIFdelend \DIFaddbegin \DIFadd{–}\DIFaddend 30,000 yen per night\DIFdelbegin \DIFdel{top 10 positive keywords list for English-speakers}\DIFdelend .

    \DIFdelbegin \DIFdel{In our }\DIFdelend \DIFaddbegin \DIFadd{For }\DIFaddend research questions \ref{rsq:hard_soft} and \ref{rsq:hard_soft_diff}, we \DIFdelbegin \DIFdel{ponder }\DIFdelend \DIFaddbegin \DIFadd{considered }\DIFaddend how customers of both cultural backgrounds \DIFdelbegin \DIFdel{evaluate }\DIFdelend \DIFaddbegin \DIFadd{evaluated the }\DIFaddend hard and soft attributes of \DIFdelbegin \DIFdel{the hotel and how they differ in those evaluations. Here}\DIFdelend \DIFaddbegin \DIFadd{hotels. Here, }\DIFaddend we define hard attributes as those \DIFdelbegin \DIFdel{relating }\DIFdelend \DIFaddbegin \DIFadd{related }\DIFaddend to the hotel's physical, structural, or environmental aspects. These are often impossible or impractical to change by the hotel management and staff, such as facilities, infrastructure, \DIFdelbegin \DIFdel{the }\DIFdelend surroundings, view, \DIFdelbegin \DIFdel{or }\DIFdelend \DIFaddbegin \DIFadd{and }\DIFaddend location convenience. \DIFdelbegin \DIFdel{On the other hand}\DIFdelend \DIFaddbegin \DIFadd{In contrast}\DIFaddend , hotel staff and management can change soft attributes, for example, by improving the hotel's services via training or hiring specialized staff, improving the quantity or quality of amenities, bedsheets, or general cleanliness. Our study \DIFdelbegin \DIFdel{found }\DIFdelend \DIFaddbegin \DIFadd{discovered }\DIFaddend that Chinese tourists \DIFdelbegin \DIFdel{are mostly positively reacting more }\DIFdelend \DIFaddbegin \DIFadd{mostly positively react }\DIFaddend to the hotel's hard attributes\DIFdelbegin \DIFdel{. There is a slightly hard leaning (53\%) concern with hard attributes in negative sentences, albeit this is }\DIFdelend \DIFaddbegin \DIFadd{, albeit the negative evaluations are }\DIFaddend more uniform than the positive evaluations\DIFaddbegin \DIFadd{, with a tendency of 53 \% towards hard attributes}\DIFaddend . English-speaking tourists, on the other hand, \DIFdelbegin \DIFdel{are both positively and negatively }\DIFdelend \DIFaddbegin \DIFadd{were }\DIFaddend more responsive to soft attributes\DIFaddbegin \DIFadd{, either positively or negatively}\DIFaddend . In the case of negative keywords, \DIFdelbegin \DIFdel{English-speaking tourists are overwhelmingly more concerned with }\DIFdelend \DIFaddbegin \DIFadd{they were more concerned about }\DIFaddend the hotel's soft attributes\DIFdelbegin \DIFdel{dissatisfaction somehow}\DIFdelend .

    One factor that both populations \DIFdelbegin \DIFdel{have }\DIFdelend \DIFaddbegin \DIFadd{had }\DIFaddend in common is \DIFaddbegin \DIFadd{that}\DIFaddend , when perceiving the hotel negatively, \DIFaddbegin \DIFadd{the }\DIFaddend ``\begin{CJK}{UTF8}{gbsn}老\end{CJK} (old)\DIFdelbegin \DIFdel{'',``dated'',``outdated'',}\DIFdelend \DIFaddbegin \DIFadd{,'' ``dated,'' ``outdated,'' }\DIFaddend or ``\begin{CJK}{UTF8}{gbsn}陈旧\end{CJK} (obsolete)'' aspects of the room or the hotel \DIFdelbegin \DIFdel{are being criticized across , surprisingly, }\DIFdelend \DIFaddbegin \DIFadd{were surprisingly criticized across }\DIFaddend most price ranges. \DIFdelbegin \DIFdel{This is, however, }\DIFdelend \DIFaddbegin \DIFadd{However, this is }\DIFaddend a hard attribute and is unlikely to change for most hotels.

  \subsection{Chinese tourists\DIFdelbegin \DIFdel{- }\DIFdelend \DIFaddbegin \DIFadd{: }\DIFaddend A big and clean space}\label{disc:zh}

    We found that mainland Chinese tourists \DIFdelbegin \DIFdel{are satisfied mostly by Japanese hotels' }\DIFdelend \DIFaddbegin \DIFadd{were mainly satisfied by }\DIFaddend big and clean spaces \DIFdelbegin \DIFdel{. From the adjectival pairings that we }\DIFdelend \DIFaddbegin \DIFadd{in Japanese hotels. The adjectival pairings }\DIFaddend extracted with dependency parsing and POS tagging \DIFdelbegin \DIFdel{in Table \ref{tab:adj_zh_pos}, we can observe that mostly they mean }\DIFdelend \DIFaddbegin \DIFadd{(Table \ref{tab:adj_zh_pos})  imply }\DIFaddend big and clean rooms. Other mentions \DIFdelbegin \DIFdel{are also }\DIFdelend \DIFaddbegin \DIFadd{included }\DIFaddend big markets nearby or a big bed. \DIFdelbegin \DIFdel{We can observe that across }\DIFdelend \DIFaddbegin \DIFadd{Across }\DIFaddend different price ranges, the usage of the word ``\begin{CJK}{UTF8}{gbsn}大\end{CJK} (big)'' \DIFdelbegin \DIFdel{increases as the hotel increases in price }\DIFdelend \DIFaddbegin \DIFadd{increased with the increasing price of the hotel}\DIFaddend . However, \DIFdelbegin \DIFdel{we can see that they still react positively in a significant manner in }\DIFdelend \DIFaddbegin \DIFadd{they still reacted positively in }\DIFaddend cheaper hotels. \DIFdelbegin \DIFdel{Inspecting }\DIFdelend \DIFaddbegin \DIFadd{When inspecting }\DIFaddend closer by taking random samples of the pairs of ``\begin{CJK}{UTF8}{gbsn}大 空间\end{CJK} (big space)'' or ``\begin{CJK}{UTF8}{gbsn}大 面积\end{CJK} (large area)\DIFdelbegin \DIFdel{'',we can see that there are }\DIFdelend \DIFaddbegin \DIFadd{,'' we notice that there were }\DIFaddend also many references to the public bathing facilities in the hotel. \DIFdelbegin \DIFdel{We can also see them mentioned as }\DIFdelend \DIFaddbegin \DIFadd{Such references were also implied by }\DIFaddend a word pairing ``\begin{CJK}{UTF8}{gbsn}棒 温泉\end{CJK} (great hot spring)\DIFdelbegin \DIFdel{''.}\DIFdelend \DIFaddbegin \DIFadd{.'' }\DIFaddend In Japan, there \DIFdelbegin \DIFdel{is what is called }\DIFdelend \DIFaddbegin \DIFadd{are the so-called }\DIFaddend ``\begin{CJK}{UTF8}{min}銭湯\end{CJK} (\textit{sent\=o})\DIFdelbegin \DIFdel{'',}\DIFdelend \DIFaddbegin \DIFadd{,'' }\DIFaddend which are artificially \DIFdelbegin \DIFdel{made }\DIFdelend \DIFaddbegin \DIFadd{constructed }\DIFaddend public bathing facilities, \DIFdelbegin \DIFdel{on occasions }\DIFdelend including saunas and baths with unique qualities. On the other hand, there are natural hot springs, called ``\begin{CJK}{UTF8}{gbsn}温泉\end{CJK} (\textit{onsen})\DIFdelbegin \DIFdel{''.They can either be }\DIFdelend \DIFaddbegin \DIFadd{.'' The word ``}\begin{CJK}{UTF8}{gbsn}\DIFadd{温泉}\end{CJK} \DIFadd{(}\textit{\DIFadd{onsen}}\DIFadd{)'' can be used }\DIFaddend bathing in the natural source of the water \DIFaddbegin \DIFadd{in rock baths }\DIFaddend or using the \DIFdelbegin \DIFdel{hot springs }\DIFdelend \DIFaddbegin \DIFadd{natural hot spring water }\DIFaddend in artificially made \DIFaddbegin \DIFadd{tiled }\DIFaddend bath facilities. It is a Japanese custom \DIFdelbegin \DIFdel{and culture }\DIFdelend that all customers \DIFdelbegin \DIFdel{use the facilities after cleaning }\DIFdelend \DIFaddbegin \DIFadd{first clean }\DIFaddend themselves in a shower and \DIFdelbegin \DIFdel{go into the baths without any clothes. It can }\DIFdelend \DIFaddbegin \DIFadd{afterward use the baths nude. It could }\DIFaddend be a cultural shock for many tourists \DIFdelbegin \DIFdel{, but this is }\DIFdelend \DIFaddbegin \DIFadd{but }\DIFaddend a fundamental attraction for many \DIFaddbegin \DIFadd{others}\DIFaddend . 

    However, the size of the room or the bed is a hard attribute. Without considering rebuilding the hotel, it is not trivial to improve \DIFdelbegin \DIFdel{on. On the other hand}\DIFdelend \DIFaddbegin \DIFadd{this attribute. In contrast}\DIFaddend , cleanliness is mostly \DIFdelbegin \DIFdel{relating }\DIFdelend \DIFaddbegin \DIFadd{related }\DIFaddend to soft attributes when we observe its adjectival pairings. We can observe pairs such as ``\begin{CJK}{UTF8}{gbsn}干净 房间\end{CJK} (clean room)'' at the top rank of all price ranges \DIFdelbegin \DIFdel{, and then variably }\DIFdelend \DIFaddbegin \DIFadd{and thereupon }\DIFaddend ``\begin{CJK}{UTF8}{gbsn}干净 酒店\end{CJK} (clean hotel)\DIFdelbegin \DIFdel{'',}\DIFdelend \DIFaddbegin \DIFadd{,'' }\DIFaddend ``\begin{CJK}{UTF8}{gbsn}干净 总体\end{CJK} (clean overall)\DIFdelbegin \DIFdel{'',}\DIFdelend \DIFaddbegin \DIFadd{,'' }\DIFaddend ``\begin{CJK}{UTF8}{gbsn}干净 环境\end{CJK} (clean environment)\DIFdelbegin \DIFdel{'',}\DIFdelend \DIFaddbegin \DIFadd{,'' }\DIFaddend and ``\begin{CJK}{UTF8}{gbsn}干净 设施\end{CJK} (clean facilities)\DIFdelbegin \DIFdel{'',}\DIFdelend \DIFaddbegin \DIFadd{,'' }\DIFaddend among other examples. In negative reviews, there \DIFdelbegin \DIFdel{is }\DIFdelend \DIFaddbegin \DIFadd{was }\DIFaddend a mention of criticizing the ``\begin{CJK}{UTF8}{gbsn}一般 卫生\end{CJK} (general hygiene)'' of the hotel, although it \DIFdelbegin \DIFdel{is }\DIFdelend \DIFaddbegin \DIFadd{was }\DIFaddend an uncommon pair. Therefore, we can assert that cleanliness \DIFdelbegin \DIFdel{is }\DIFdelend \DIFaddbegin \DIFadd{was }\DIFaddend an important soft attribute for Chinese customers and \DIFdelbegin \DIFdel{that they are mostly pleased with their expectations being met}\DIFdelend \DIFaddbegin \DIFadd{they were mostly pleased when their expectations were fulfilled}\DIFaddend . 

    \DIFdelbegin \DIFdel{One key component we }\DIFdelend \DIFaddbegin \DIFadd{A key component }\DIFaddend found in Chinese customer satisfaction soft factors \DIFdelbegin \DIFdel{is }\DIFdelend \DIFaddbegin \DIFadd{was }\DIFaddend the inclusion of breakfast within the hotel. While other food-related words were extracted, most of them were general, \DIFdelbegin \DIFdel{like }\DIFdelend \DIFaddbegin \DIFadd{such as }\DIFaddend ``food'' or ``eating,'' \DIFdelbegin \DIFdel{which were lower ranking}\DIFdelend \DIFaddbegin \DIFadd{and were lower-ranking}\DIFaddend . In contrast, the word ``\begin{CJK}{UTF8}{gbsn}早餐\end{CJK} (breakfast)\DIFdelbegin \DIFdel{'' refers possibly }\DIFdelend \DIFaddbegin \DIFadd{,'' which possibly referred }\DIFaddend to its inclusion in the hotel commodities, was frequently used in positive texts compared to other food-related words. The word ``\begin{CJK}{UTF8}{gbsn}早餐\end{CJK} (breakfast)'' \DIFdelbegin \DIFdel{is also observed across }\DIFdelend \DIFaddbegin \DIFadd{was also present in }\DIFaddend all price ranges, \DIFdelbegin \DIFdel{although }\DIFdelend \DIFaddbegin \DIFadd{albeit }\DIFaddend at different priorities in each of them. \DIFdelbegin \DIFdel{However, we assert that it is }\DIFdelend \DIFaddbegin \DIFadd{For this reason, we regard it as }\DIFaddend an important factor. \DIFdelbegin \DIFdel{Observing word pairs from }\DIFdelend \DIFaddbegin \DIFadd{From the word pairs of }\DIFaddend the positive Chinese keywords in Table \ref{tab:adj_zh_pos}, we can also \DIFdelbegin \DIFdel{see }\DIFdelend \DIFaddbegin \DIFadd{note }\DIFaddend that ``\begin{CJK}{UTF8}{gbsn}不错\end{CJK} (not bad)'' is paired \DIFdelbegin \DIFdel{as }\DIFdelend \DIFaddbegin \DIFadd{with }\DIFaddend ``\begin{CJK}{UTF8}{gbsn}不错 早餐\end{CJK} (nice breakfast)'' in four of the seven price ranges with reviews available as part of the top \DIFdelbegin \DIFdel{4 }\DIFdelend \DIFaddbegin \DIFadd{four }\DIFaddend pairings. It is only slightly lower \DIFdelbegin \DIFdel{on }\DIFdelend \DIFaddbegin \DIFadd{in }\DIFaddend other categories, although it is not \DIFdelbegin \DIFdel{shown }\DIFdelend \DIFaddbegin \DIFadd{depicted }\DIFaddend on the table. Thus\DIFaddbegin \DIFadd{, }\DIFaddend we consider that a recommended strategy for hotel management is to invest in the inclusion or \DIFdelbegin \DIFdel{betterment }\DIFdelend \DIFaddbegin \DIFadd{improvement }\DIFaddend of hotel breakfast to increase \DIFaddbegin \DIFadd{the number of }\DIFaddend good reviews.

  \subsection{Western tourists\DIFdelbegin \DIFdel{- }\DIFdelend \DIFaddbegin \DIFadd{: }\DIFaddend A friendly face \DIFdelbegin \DIFdel{, }\DIFdelend and absolutely clean}\label{disc:en}

    From the satisfaction factors of English-speaking tourists, we can see that at least three words \DIFdelbegin \DIFdel{relate directly }\DIFdelend \DIFaddbegin \DIFadd{were directly related }\DIFaddend to staff friendliness and services \DIFdelbegin \DIFdel{, being ``staff'',}\DIFdelend \DIFaddbegin \DIFadd{in the general database, }\DIFaddend ``\DIFdelbegin \DIFdel{helpful}\DIFdelend \DIFaddbegin \DIFadd{staff,'' ``helpful,}\DIFaddend '' and ``friendliness\DIFdelbegin \DIFdel{'' in the general database.}\DIFdelend \DIFaddbegin \DIFadd{.'' }\DIFaddend The word ``staff'' is the second most frequently used word for satisfied customers across most price ranges \DIFdelbegin \DIFdel{, }\DIFdelend and only third in one of them. \DIFdelbegin \DIFdel{Adding to that, }\DIFdelend \DIFaddbegin \DIFadd{In addition, the words }\DIFaddend ``helpful'' and ``friendly'' follow it lower in the list in most price ranges. The word ``good'' \DIFdelbegin \DIFdel{is mostly about }\DIFdelend \DIFaddbegin \DIFadd{mainly refers to }\DIFaddend the location, \DIFdelbegin \DIFdel{the }\DIFdelend service, breakfast, or English availability in Table \ref{tab:adj_en_pos}. \DIFdelbegin \DIFdel{Like }\DIFdelend \DIFaddbegin \DIFadd{Similar to }\DIFaddend Chinese customers, Western customers also \DIFdelbegin \DIFdel{seem }\DIFdelend \DIFaddbegin \DIFadd{seemed }\DIFaddend to enjoy the included breakfasts regarding their satisfaction keyword pairings. However, the \DIFaddbegin \DIFadd{relevant }\DIFaddend word does not appear \DIFdelbegin \DIFdel{directly }\DIFdelend in the top 10 list \DIFdelbegin \DIFdel{as in }\DIFdelend \DIFaddbegin \DIFadd{directly, in contrast to }\DIFaddend their Chinese counterparts. The words ``helpful'' and ``friendly'' are mostly paired with ``staff\DIFdelbegin \DIFdel{'',``concierge'',``desk'',}\DIFdelend \DIFaddbegin \DIFadd{,'' ``concierge,'' ``desk,'' }\DIFaddend and ``service\DIFdelbegin \DIFdel{''.When we look at }\DIFdelend \DIFaddbegin \DIFadd{.'' By considering }\DIFaddend the negative keyword \DIFdelbegin \DIFdel{‘poor’ }\DIFdelend \DIFaddbegin \DIFadd{``poor'' }\DIFaddend and its pairings in Table \ref{tab:adj_en_neg}, we \DIFdelbegin \DIFdel{see that it is also }\DIFdelend \DIFaddbegin \DIFadd{realized once again that Western tourists were disappointed with }\DIFaddend service-related \DIFdelbegin \DIFdel{concepts that the Western tourists are disappointed with when they react }\DIFdelend \DIFaddbegin \DIFadd{concepts and reacted }\DIFaddend negatively.

    Another soft attribute that is high on the list for most of the price ranges is the word ``clean''\DIFdelbegin \DIFdel{. Since it }\DIFdelend \DIFaddbegin \DIFadd{; because this }\DIFaddend is an adjective, we \DIFdelbegin \DIFdel{have }\DIFdelend \DIFaddbegin \DIFadd{also }\DIFaddend explored the word pairings\DIFdelbegin \DIFdel{as well. Customers are mostly praising }\DIFdelend \DIFaddbegin \DIFadd{. Customers largely praised }\DIFaddend ``clean rooms'' and ``clean bathrooms'' \DIFdelbegin \DIFdel{, while also referring }\DIFdelend \DIFaddbegin \DIFadd{and also referred }\DIFaddend to the hotel in general. \DIFdelbegin \DIFdel{It seems that when }\DIFdelend \DIFaddbegin \DIFadd{When }\DIFaddend observing the negative keyword frequencies for \DIFdelbegin \DIFdel{English-speakers}\DIFdelend \DIFaddbegin \DIFadd{English speakers}\DIFaddend , we can find words such as ``dirty'' \DIFdelbegin \DIFdel{, }\DIFdelend \DIFaddbegin \DIFadd{and }\DIFaddend ``carpet'' \DIFdelbegin \DIFdel{, and from the word pairings }\DIFdelend \DIFaddbegin \DIFadd{as well as word pairings such as }\DIFaddend ``dirty carpet\DIFdelbegin \DIFdel{'',}\DIFdelend \DIFaddbegin \DIFadd{,'' }\DIFaddend ``dirty room\DIFdelbegin \DIFdel{'',}\DIFdelend \DIFaddbegin \DIFadd{,'' }\DIFaddend and ``dirty bathroom\DIFdelbegin \DIFdel{''.}\DIFdelend \DIFaddbegin \DIFadd{.'' }\DIFaddend Along with complaints about off-putting smells, we \DIFdelbegin \DIFdel{can }\DIFdelend \DIFaddbegin \DIFadd{could }\DIFaddend conclude that Western tourists \DIFdelbegin \DIFdel{have }\DIFdelend \DIFaddbegin \DIFadd{had }\DIFaddend high expectations about cleanliness when traveling in Japan.

    An interesting detail of the keyword ranking is that the word ``comfortable'' \DIFdelbegin \DIFdel{is }\DIFdelend \DIFaddbegin \DIFadd{was }\DIFaddend high on the satisfaction factors and ``uncomfortable'' \DIFaddbegin \DIFadd{was }\DIFaddend high on the dissatisfaction factors. The words \DIFdelbegin \DIFdel{are }\DIFdelend \DIFaddbegin \DIFadd{were }\DIFaddend paired with nouns \DIFdelbegin \DIFdel{like ``bed'',or ``room'',``pillow'' or ``mattress'',generally referring }\DIFdelend \DIFaddbegin \DIFadd{such as ``bed,'' ``room,'' ``pillow,'' and ``mattress,'' when they generally referred }\DIFaddend to their sleep conditions in the hotel.
    It seems that Western tourists \DIFdelbegin \DIFdel{are highly sensitive to comfort levels in the hotelsand whether it reaches }\DIFdelend \DIFaddbegin \DIFadd{were particularly sensitive about the hotels’ comfort levels and whether they reached }\DIFaddend their expectations. The ranking for the negative keyword ``uncomfortable'' is similar across most price ranges \DIFdelbegin \DIFdel{, }\DIFdelend except the two most expensive ones, where this keyword disappears from the top 10 list.

    \DIFdelbegin \DIFdel{While less high }\DIFdelend \DIFaddbegin \DIFadd{Albeit lower }\DIFaddend in priority, the price range of 15,000 to 20,000 yen hotels also \DIFdelbegin \DIFdel{mentions }\DIFdelend \DIFaddbegin \DIFadd{includes }\DIFaddend ``free'' as one of the top 10 positive keywords, \DIFdelbegin \DIFdel{paired mostly with ``wifi''.}\DIFdelend \DIFaddbegin \DIFadd{mainly paired with ``Wi-Fi.'' }\DIFaddend This price range \DIFdelbegin \DIFdel{is mostly for }\DIFdelend \DIFaddbegin \DIFadd{corresponds to }\DIFaddend business hotels, where \DIFdelbegin \DIFdel{we infer users would be expecting }\DIFdelend \DIFaddbegin \DIFadd{users would expect }\DIFaddend this feature the most. 
\DIFdelbegin \DIFdel{Western tourists are highly sensitive to comfort levels in the hotels and whether it reaches their expectations.
}\DIFdelend 


  \subsection{Tobacco, what \DIFdelbegin \DIFdel{'s }\DIFdelend \DIFaddbegin \DIFadd{is }\DIFaddend that smell?}\label{disc:tobacco}

    A concern for Western tourists was the smell of tobacco in their room, which can be \DIFdelbegin \DIFdel{considered }\DIFdelend \DIFaddbegin \DIFadd{regarded as }\DIFaddend a soft attribute. Tobacco was found not only as a standalone word \DIFaddbegin \DIFadd{with }\DIFaddend ``cigarette'' \DIFdelbegin \DIFdel{, but also as }\DIFdelend \DIFaddbegin \DIFadd{but also in }\DIFaddend word pairs in Table \ref{tab:adj_en_neg}\DIFdelbegin \DIFdel{. We can find other related word pairs }\DIFdelend \DIFaddbegin \DIFadd{, }\DIFaddend such as ``funny smell\DIFdelbegin \DIFdel{''.Upon manual inspection of }\DIFdelend \DIFaddbegin \DIFadd{.'' By manually inspecting }\DIFaddend a sample of reviews with this keyword, we \DIFdelbegin \DIFdel{found }\DIFdelend \DIFaddbegin \DIFadd{noticed }\DIFaddend that the room was often advertised as non-smoking\DIFdelbegin \DIFdel{, yet, }\DIFdelend \DIFaddbegin \DIFadd{; however, }\DIFaddend the smell permeated the room and curtains. Another common complaint was that there were no \DIFdelbegin \DIFdel{non-smoking facilities availableat all in the first place}\DIFdelend \DIFaddbegin \DIFadd{nonsmoking facilities available}\DIFaddend . The smell of smoke can completely ruin some customers’ stay and \DIFdelbegin \DIFdel{give a bad impression to review writers, }\DIFdelend \DIFaddbegin \DIFadd{thus lead to bad reviews, thereby }\DIFaddend lowering the number of future customers.

    \DIFdelbegin \DIFdel{However, in comparison, Chinese customers seem not to }\DIFdelend \DIFaddbegin \DIFadd{In contrast, Chinese customers seemed to not }\DIFaddend be bothered by this\DIFdelbegin \DIFdel{at all. We consulted studies involving the use of tobacco in different countries. Previous research states }\DIFdelend \DIFaddbegin \DIFadd{. Previous research has stated }\DIFaddend that 49\DIFdelbegin \DIFdel{- }\DIFdelend \DIFaddbegin \DIFadd{–}\DIFaddend 60 \% of Chinese men (and 2.0\DIFdelbegin \DIFdel{- }\DIFdelend \DIFaddbegin \DIFadd{–}\DIFaddend 2.8 \% of women) currently smoke or \DIFdelbegin \DIFdel{have smoked before}\DIFdelend \DIFaddbegin \DIFadd{smoked in the past}\DIFaddend . This was \DIFdelbegin \DIFdel{taken }\DIFdelend \DIFaddbegin \DIFadd{derived }\DIFaddend from a sample of 170,000 Chinese adults in \DIFdelbegin \DIFdel{2013-2014}\DIFdelend \DIFaddbegin \DIFadd{2013–2014}\DIFaddend , which is high compared to many English-speaking countries \cite[][]{zhang2019tobacco,who2015tobacco}.

    Japan has a polarized view on the topic of smoking. \DIFdelbegin \DIFdel{Despite being }\DIFdelend \DIFaddbegin \DIFadd{Although it has }\DIFaddend one of the world’s largest tobacco markets, \DIFdelbegin \DIFdel{its }\DIFdelend \DIFaddbegin \DIFadd{tobacco }\DIFaddend use has decreased in recent years. Smoking in public spaces is prohibited in some wards of Tokyo (namely Chiyoda, Shinjuku, and Shibuya). However, it is generally only \DIFdelbegin \DIFdel{urged }\DIFdelend \DIFaddbegin \DIFadd{suggested }\DIFaddend and not mandatory to \DIFdelbegin \DIFdel{have }\DIFdelend \DIFaddbegin \DIFadd{lift }\DIFaddend smoking restrictions in restaurants, bars, hotels, and public areas. However, many places have designated smoking rooms \DIFdelbegin \DIFdel{are available }\DIFdelend to keep the smoke in an enclosed area and avoid bothering others. \DIFdelbegin \DIFdel{Despite this}\DIFdelend \DIFaddbegin \DIFadd{Nevertheless}\DIFaddend , businesses, especially those who cater to certain customers, \DIFdelbegin \DIFdel{will generally be discouraged from having }\DIFdelend \DIFaddbegin \DIFadd{are generally discouraged by }\DIFaddend smoking restrictions if they want to \DIFdelbegin \DIFdel{keep }\DIFdelend \DIFaddbegin \DIFadd{maintain }\DIFaddend their clientele. \DIFdelbegin \DIFdel{If Japanese hotels want to }\DIFdelend \DIFaddbegin \DIFadd{To }\DIFaddend cater to all kinds of customers, \DIFaddbegin \DIFadd{including }\DIFaddend Western and Asian\DIFdelbegin \DIFdel{alike, they }\DIFdelend \DIFaddbegin \DIFadd{, Japanese hotels }\DIFaddend must provide spaces without tobacco smell. \DIFdelbegin \DIFdel{After all, even if it }\DIFdelend \DIFaddbegin \DIFadd{Even if the smoke }\DIFaddend does not bother a few customers, the lack of \DIFaddbegin \DIFadd{such a }\DIFaddend smell would make it an appropriate space for all customers.


  \subsection{Location, location, location}\label{disc:location}

    The hotel's location, closeness to the subway and public \DIFdelbegin \DIFdel{transport, and nearby shops' availability were observed }\DIFdelend \DIFaddbegin \DIFadd{transportation, and availability of nearby shops proved }\DIFaddend to be of importance to both Chinese and English-speaking tourists. In positive word pairings \DIFaddbegin \DIFadd{in }\DIFaddend Tables \ref{tab:adj_zh_pos} and \ref{tab:adj_en_pos}, we can find pairs such as ``\begin{CJK}{UTF8}{gbsn}不错 位置\end{CJK} (nice location)\DIFdelbegin \DIFdel{'',}\DIFdelend \DIFaddbegin \DIFadd{,'' }\DIFaddend ``\begin{CJK}{UTF8}{gbsn}近 地铁站\end{CJK} (near subway station)\DIFdelbegin \DIFdel{'',}\DIFdelend \DIFaddbegin \DIFadd{,'' }\DIFaddend ``\begin{CJK}{UTF8}{gbsn}近 地铁\end{CJK} (near subway)'' in Chinese texts and ``good location\DIFdelbegin \DIFdel{'',}\DIFdelend \DIFaddbegin \DIFadd{,'' }\DIFaddend ``great location\DIFdelbegin \DIFdel{'',}\DIFdelend \DIFaddbegin \DIFadd{,'' }\DIFaddend and ``great view'' \DIFdelbegin \DIFdel{, }\DIFdelend as well as single keywords ``location'' and ``shopping'' for \DIFdelbegin \DIFdel{English-speakers}\DIFdelend \DIFaddbegin \DIFadd{English speakers}\DIFaddend , and ``\begin{CJK}{UTF8}{gbsn}交通\end{CJK} (traffic)\DIFdelbegin \DIFdel{'',}\DIFdelend \DIFaddbegin \DIFadd{,'' }\DIFaddend ``\begin{CJK}{UTF8}{gbsn}购物\end{CJK} (shopping)\DIFdelbegin \DIFdel{'',}\DIFdelend \DIFaddbegin \DIFadd{,'' }\DIFaddend ``\begin{CJK}{UTF8}{gbsn}地铁\end{CJK} (subway)\DIFdelbegin \DIFdel{'',}\DIFdelend \DIFaddbegin \DIFadd{,'' }\DIFaddend and ``\begin{CJK}{UTF8}{gbsn}环境\end{CJK} (environment or surroundings)'' for Chinese speakers. All of these keywords and their location in each population's priorities across the price ranges \DIFdelbegin \DIFdel{signal that while it was not the priority for either of them, the }\DIFdelend \DIFaddbegin \DIFadd{signify that the }\DIFaddend hotel's location \DIFdelbegin \DIFdel{is }\DIFdelend \DIFaddbegin \DIFadd{was }\DIFaddend a secondary but still important point \DIFdelbegin \DIFdel{in the hotel's }\DIFdelend \DIFaddbegin \DIFadd{for their }\DIFaddend satisfaction. However, since this is a hard attribute, \DIFdelbegin \DIFdel{unchangeable to the hotel's management, }\DIFdelend it is not often considered in the literature. \DIFdelbegin \DIFdel{Upon inspection of }\DIFdelend \DIFaddbegin \DIFadd{By examining }\DIFaddend examples from the data, we \DIFdelbegin \DIFdel{found }\DIFdelend \DIFaddbegin \DIFadd{recognized }\DIFaddend that most customers were satisfied if the hotel was near \DIFdelbegin \DIFdel{to }\DIFdelend at least two \DIFdelbegin \DIFdel{other subjects}\DIFdelend \DIFaddbegin \DIFadd{of the following facilities}\DIFaddend : subway, train, and convenience stores. 

    Japan is a country with a peculiar public \DIFdelbegin \DIFdel{transport system. The rush hourmakes for a subway filled to the brim with people in suits making their commute}\DIFdelend \DIFaddbegin \DIFadd{transportation system. During rush hour, the subway is crowded with commuters}\DIFaddend , and trains and subway stations \DIFdelbegin \DIFdel{in Tokyo }\DIFdelend create a confusing public \DIFdelbegin \DIFdel{transport }\DIFdelend \DIFaddbegin \DIFadd{transportation }\DIFaddend map for a visitor \DIFaddbegin \DIFadd{in Tokyo}\DIFaddend . Buses are also available, \DIFdelbegin \DIFdel{although }\DIFdelend \DIFaddbegin \DIFadd{albeit }\DIFaddend less used than \DIFdelbegin \DIFdel{the }\DIFdelend rail systems in \DIFdelbegin \DIFdel{the big }\DIFdelend \DIFaddbegin \DIFadd{metropolitan }\DIFaddend cities. These three \DIFdelbegin \DIFdel{are unusually }\DIFdelend \DIFaddbegin \DIFadd{means of transportation are usually }\DIFaddend affordable in price. \DIFdelbegin \DIFdel{Then there are the more expensive transports}\DIFdelend \DIFaddbegin \DIFadd{There are more expensive means}\DIFaddend , such as the bullet train \textit{shinkansen} for traveling across the country \DIFdelbegin \DIFdel{, }\DIFdelend and taxis. \DIFdelbegin \DIFdel{Taxis in Japan are a luxury }\DIFdelend \DIFaddbegin \DIFadd{The latter is a luxury in Japan }\DIFaddend compared to other countries\DIFdelbegin \DIFdel{. In }\DIFdelend \DIFaddbegin \DIFadd{, especially in }\DIFaddend less developed countries, \DIFdelbegin \DIFdel{a taxi is the cheap method of transport of }\DIFdelend \DIFaddbegin \DIFadd{where this is a low-cost }\DIFaddend choice. In Japan, taxis \DIFdelbegin \DIFdel{are made to }\DIFdelend provide a high-quality experience with a matching price. \DIFdelbegin \DIFdel{This means that for tourists}\DIFdelend \DIFaddbegin \DIFadd{Therefore}\DIFaddend , subway availability and maps or GPS applications \DIFdelbegin \DIFdel{and }\DIFdelend \DIFaddbegin \DIFadd{as well as }\DIFaddend a plan to travel the city are of utmost necessity \DIFaddbegin \DIFadd{for tourists}\DIFaddend . 

    Japanese convenience stores \DIFdelbegin \DIFdel{, on the other hand, }\DIFdelend are also famous worldwide\DIFdelbegin \DIFdel{. Japanese convenience stores are a haven for the traveler in need. It offers anything}\DIFdelend \DIFaddbegin \DIFadd{, because they offer a wide range of services and products}\DIFaddend , from drinks and snacks to full meals, copy and scanning machines, alcohol, cleaning supplies, personal hygiene items, underwear, towels, \DIFdelbegin \DIFdel{international ATMs, among other things}\DIFdelend \DIFaddbegin \DIFadd{and international ATMs}\DIFaddend . If some trouble \DIFdelbegin \DIFdel{occurred}\DIFdelend \DIFaddbegin \DIFadd{occurs}\DIFaddend , or a traveler forgot to pack a particular item, it is \DIFdelbegin \DIFdel{almost sure }\DIFdelend \DIFaddbegin \DIFadd{mostly certain }\DIFaddend that they can find it. 

    Therefore, considering that both \DIFdelbegin \DIFdel{transport }\DIFdelend \DIFaddbegin \DIFadd{transportation }\DIFaddend systems and nearby shops are points of interest for Chinese and Western tourists, \DIFdelbegin \DIFdel{Japanese hotels have to carefully choose their location from the moment they are constructed. While }\DIFdelend \DIFaddbegin \DIFadd{the location of a Japanese hotel must be carefully chosen prior to construction. Although }\DIFaddend not a top priority, this is a universal factor for both customer groups \DIFdelbegin \DIFdel{, and it can be an instant way }\DIFdelend \DIFaddbegin \DIFadd{and is conducive }\DIFaddend to generate positive reviews.

\section{Discussion}\label{discussion}

  \DIFdelbegin \DIFdel{Below }\DIFdelend \DIFaddbegin \DIFadd{In this section, }\DIFaddend we explore the possible interactions with Japanese hospitality, the differences between \DIFaddbegin \DIFadd{perceptions of }\DIFaddend Chinese and Western tourists \DIFdelbegin \DIFdel{, the possible cause for them, how they vary }\DIFdelend \DIFaddbegin \DIFadd{as well as the possible causes, variations in perceptions }\DIFaddend across different price ranges, and \DIFdelbegin \DIFdel{what they imply }\DIFdelend \DIFaddbegin \DIFadd{the implications }\DIFaddend for the industry. We also discuss the differences between the hotel's hard and soft attributes and \DIFdelbegin \DIFdel{how they contribute }\DIFdelend \DIFaddbegin \DIFadd{their contributions }\DIFaddend to customers' satisfaction.

  \subsection{Western and Chinese tourists in the Japanese hospitality environment}\label{disc:omotenashi}

    To \DIFdelbegin \DIFdel{this day, scholars continue to correct }\DIFdelend \DIFaddbegin \DIFadd{date, scholars have been correcting }\DIFaddend our historical bias towards the \DIFdelbegin \DIFdel{west}\DIFdelend \DIFaddbegin \DIFadd{West}\DIFaddend . Studies have determined that different cultural backgrounds lead to different expectations, which \DIFdelbegin \DIFdel{influences }\DIFdelend \DIFaddbegin \DIFadd{influence }\DIFaddend tourists' satisfaction. \DIFdelbegin \DIFdel{Meaning}\DIFdelend \DIFaddbegin \DIFadd{In other words}\DIFaddend , tourists of a particular culture \DIFdelbegin \DIFdel{will }\DIFdelend have different leading satisfaction factors across different destinations. However, Japan presents a particular environment\DIFdelbegin \DIFdel{. The }\DIFdelend \DIFaddbegin \DIFadd{; the }\DIFaddend spirit of hospitality and service, \textit{omotenashi}, excels and is considered \DIFaddbegin \DIFadd{to be of }\DIFaddend the highest standard across the world. \DIFdelbegin \DIFdel{Can }\DIFdelend \DIFaddbegin \DIFadd{Our study explores whether }\DIFaddend such an environment \DIFaddbegin \DIFadd{can }\DIFaddend affect different cultures equally \DIFdelbegin \DIFdel{? Or is it only attractive }\DIFdelend \DIFaddbegin \DIFadd{Or whether it is attractive only }\DIFaddend to certain cultures\DIFdelbegin \DIFdel{? Our study brings light to these questions}\DIFdelend .

    Our results indicate that \DIFdelbegin \DIFdel{out of the two, }\DIFdelend Western tourists are \DIFdelbegin \DIFdel{the most }\DIFdelend \DIFaddbegin \DIFadd{more }\DIFaddend satisfied with soft attributes, such as friendly and helpful staff in Japan\DIFaddbegin \DIFadd{, than Chinese tourists}\DIFaddend . As explained earlier in this paper, Japan is \DIFdelbegin \DIFdel{famous }\DIFdelend \DIFaddbegin \DIFadd{well known }\DIFaddend for its customer service\DIFdelbegin \DIFdel{all over the world}\DIFdelend . Respectful language and bowing are not exclusive to \DIFdelbegin \DIFdel{high priced }\DIFdelend \DIFaddbegin \DIFadd{high-priced }\DIFaddend hotels or businesses\DIFdelbegin \DIFdel{. These can even be found }\DIFdelend \DIFaddbegin \DIFadd{; these are met }\DIFaddend in convenience stores \DIFaddbegin \DIFadd{as well}\DIFaddend . The level of hospitality\DIFdelbegin \DIFdel{in even the cheapest of convenience stores }\DIFdelend \DIFaddbegin \DIFadd{, even in the cheapest convenience store, }\DIFaddend is starkly different from \DIFdelbegin \DIFdel{Westerner experiences. While }\DIFdelend \DIFaddbegin \DIFadd{Western culture. Although }\DIFaddend it could be a \DIFdelbegin \DIFdel{culture }\DIFdelend \DIFaddbegin \DIFadd{cultural }\DIFaddend shock to some, it is mostly \DIFdelbegin \DIFdel{seen positively. After all, the Japanese staff respectfully }\DIFdelend \DIFaddbegin \DIFadd{approved. The Japanese staff }\DIFaddend treats all customers \DIFdelbegin \DIFdel{. However, for some customers, this could be the best way they have been treated until that moment. Now, in higher priced }\DIFdelend \DIFaddbegin \DIFadd{respectfully, and this might be unforeseen by some customers. In higher-priced }\DIFaddend hotels, the adjectives used to praise the service \DIFdelbegin \DIFdel{go }\DIFdelend \DIFaddbegin \DIFadd{ranged }\DIFaddend from normal descriptors like ``good'' to higher levels of praise like ``wonderful staff\DIFdelbegin \DIFdel{'',}\DIFdelend \DIFaddbegin \DIFadd{,'' }\DIFaddend ``wonderful experience\DIFdelbegin \DIFdel{'',}\DIFdelend \DIFaddbegin \DIFadd{,'' }\DIFaddend ``excellent service\DIFdelbegin \DIFdel{'',}\DIFdelend \DIFaddbegin \DIFadd{,'' }\DIFaddend and ``excellent staff\DIFdelbegin \DIFdel{''.We can also see that }\DIFdelend \DIFaddbegin \DIFadd{.'' Furthermore, }\DIFaddend \cite{kozak2002} and \cite{shanka2004} have also \DIFdelbegin \DIFdel{found }\DIFdelend \DIFaddbegin \DIFadd{proven }\DIFaddend that hospitality and staff friendliness \DIFdelbegin \DIFdel{is a vital determinant in }\DIFdelend \DIFaddbegin \DIFadd{are two determinants of }\DIFaddend Western tourists' satisfaction.

    However, \DIFdelbegin \DIFdel{we can see from }\DIFdelend the negative English keywords \DIFdelbegin \DIFdel{that a big }\DIFdelend \DIFaddbegin \DIFadd{indicate that a large }\DIFaddend part of the dissatisfaction with Japanese hotels \DIFdelbegin \DIFdel{stems }\DIFdelend \DIFaddbegin \DIFadd{stemmed }\DIFaddend from a lack of hygiene and room cleanliness. Although Chinese customers \DIFdelbegin \DIFdel{only had }\DIFdelend \DIFaddbegin \DIFadd{had solely }\DIFaddend positive keywords about cleanliness, English-speaking customers \DIFdelbegin \DIFdel{have found }\DIFdelend \DIFaddbegin \DIFadd{deemed }\DIFaddend many places unacceptable to their standards\DIFdelbegin \DIFdel{. This is particularly true at hotels }\DIFdelend \DIFaddbegin \DIFadd{, particularly hotels with rooms priced }\DIFaddend below 50,000 yen per night. The most common complaint regarding cleanliness was about the carpet, followed by complaints about cigarette \DIFdelbegin \DIFdel{stench and general dirtiness}\DIFdelend \DIFaddbegin \DIFadd{smell and lack of general hygiene}\DIFaddend . \cite{kozak2002} also \DIFdelbegin \DIFdel{found }\DIFdelend \DIFaddbegin \DIFadd{proved }\DIFaddend that hygiene and cleanliness were essential satisfaction determinants for Western tourists. However, in the previous literature, this was linked merely to satisfaction. In \DIFdelbegin \DIFdel{comparison}\DIFdelend \DIFaddbegin \DIFadd{contrast}\DIFaddend , our research \DIFdelbegin \DIFdel{uncovered that words relating to cleanliness are }\DIFdelend \DIFaddbegin \DIFadd{revealed that words related to cleanliness were }\DIFaddend mostly linked to dissatisfaction. \DIFdelbegin \DIFdel{Westerners could be said to have }\DIFdelend \DIFaddbegin \DIFadd{We could assert that Westerners had }\DIFaddend a high standard of room cleanliness \DIFdelbegin \DIFdel{when }\DIFdelend compared to their Chinese counterparts.

    According to previous research, \DIFdelbegin \DIFdel{we can see that }\DIFdelend Western tourists are already inclined to appreciate hospitality for their satisfaction. When presented with Japanese hospitality, this expectation is met and overcome. In contrast, \DIFdelbegin \DIFdel{we can see from our resultsthat Chinese tourists had less focus on }\DIFdelend \DIFaddbegin \DIFadd{according to our results, Chinese tourists were more concerned about room quality rather than }\DIFaddend hospitality, staff, or service\DIFdelbegin \DIFdel{and were more concerned with room quality}\DIFdelend . However, when analyzing the word pairs for ``\begin{CJK}{UTF8}{gbsn}不错\end{CJK} (not bad)'' and \DIFdelbegin \DIFdel{for }\DIFdelend ``\begin{CJK}{UTF8}{gbsn}棒\end{CJK} (great)\DIFdelbegin \DIFdel{'',}\DIFdelend \DIFaddbegin \DIFadd{,'' }\DIFaddend we can see that they \DIFdelbegin \DIFdel{do }\DIFdelend praise staff, service\DIFaddbegin \DIFadd{, }\DIFaddend and breakfast. \DIFdelbegin \DIFdel{Observing }\DIFdelend \DIFaddbegin \DIFadd{By observing }\DIFaddend the percentage of hard to soft attributes in Figure \ref{fig:hard_soft_zh}, however, we \DIFdelbegin \DIFdel{know }\DIFdelend \DIFaddbegin \DIFadd{discover }\DIFaddend that Chinese customers \DIFdelbegin \DIFdel{are satisfied more }\DIFdelend \DIFaddbegin \DIFadd{were more satisfied }\DIFaddend with hard attributes \DIFdelbegin \DIFdel{, compared to the Western touristswho seem }\DIFdelend \DIFaddbegin \DIFadd{compared }\DIFaddend to \DIFaddbegin \DIFadd{Western tourists, who seemed to }\DIFaddend be meeting more than their expectations.

    It could be \DIFaddbegin \DIFadd{considered }\DIFaddend that Chinese culture does not expect high-level service initially. When an expectation that is not held is met, the satisfaction \DIFdelbegin \DIFdel{that stems from this }\DIFdelend \DIFaddbegin \DIFadd{derived }\DIFaddend is less than \DIFaddbegin \DIFadd{that }\DIFaddend if it was expected. \DIFdelbegin \DIFdel{On the other hand, we have the phenomenon of }\DIFdelend \DIFaddbegin \DIFadd{In contrast, some tourists report }\DIFaddend a ``nice surprise'': \DIFdelbegin \DIFdel{When }\DIFdelend \DIFaddbegin \DIFadd{when }\DIFaddend an unknown need is unexpectedly met, there is more satisfaction. It is necessary to note the difference between these two \DIFdelbegin \DIFdel{phenomenons}\DIFdelend \DIFaddbegin \DIFadd{reactions}\DIFaddend . The ``nice surprise'' \DIFaddbegin \DIFadd{reaction }\DIFaddend fulfills a need unexpectedly. Perhaps the hospitality grade in Japan does not fulfill a \DIFdelbegin \DIFdel{high enough need }\DIFdelend \DIFaddbegin \DIFadd{need high enough }\DIFaddend for the Chinese population, \DIFaddbegin \DIFadd{thereby }\DIFaddend resulting in less satisfaction. For greater satisfaction, \DIFdelbegin \DIFdel{the existence of a need being metis necessary}\DIFdelend \DIFaddbegin \DIFadd{a need must be met}\DIFaddend . However, the word ``not bad'' is at the top of the list \DIFdelbegin \DIFdel{at }\DIFdelend \DIFaddbegin \DIFadd{in }\DIFaddend most price ranges, and one of the uses is related to service. Thus, we cannot \DIFdelbegin \DIFdel{say that they are }\DIFdelend \DIFaddbegin \DIFadd{conclude that they were }\DIFaddend not satisfied with \DIFdelbegin \DIFdel{this matter. Rather, they hold }\DIFdelend \DIFaddbegin \DIFadd{the service. Instead, they held }\DIFaddend other factors at a higher priority\DIFdelbegin \DIFdel{, considering }\DIFdelend \DIFaddbegin \DIFadd{; thus, }\DIFaddend the keyword frequency \DIFdelbegin \DIFdel{is }\DIFdelend \DIFaddbegin \DIFadd{was }\DIFaddend higher for other pairings.

    Another possibility \DIFdelbegin \DIFdel{presents itself }\DIFdelend \DIFaddbegin \DIFadd{occurs }\DIFaddend when we observe the Chinese tourists’ dissatisfaction factors. Chinese tourists may have expectations about the Chinese visitors' treatment that are not being met, even in this \DIFdelbegin \DIFdel{high standard }\DIFdelend \DIFaddbegin \DIFadd{high-standard }\DIFaddend hospitality environment. \DIFdelbegin \DIFdel{Japan is known worldwide for their hospitality, but they are also known historically to be monolingual and have }\DIFdelend \DIFaddbegin \DIFadd{This could be because Japan is monolingual and has }\DIFaddend a relatively large language barrier \DIFaddbegin \DIFadd{to tourists }\DIFaddend \cite[][]{heinrich2012making,coulmas2002japan}. While the Japanese effort to accommodate English speakers is slowly \DIFdelbegin \DIFdel{taking shape, }\DIFdelend \DIFaddbegin \DIFadd{developing, efforts for }\DIFaddend Chinese accommodations can be lagging. Chinese language pamphlets \DIFdelbegin \DIFdel{, }\DIFdelend \DIFaddbegin \DIFadd{as well as }\DIFaddend Chinese texts on instructions for the hotel room \DIFdelbegin \DIFdel{, }\DIFdelend and its appliances and features (e.g., T.V. channels, Wi-Fi setup\DIFaddbegin \DIFadd{, etc.}\DIFaddend ), or \DIFdelbegin \DIFdel{just }\DIFdelend the treatment towards Chinese people could be examples. \DIFdelbegin \DIFdel{It is natural to be dissatisfied since traveling in a strange land without knowing the language can be a daunting experience}\DIFdelend \DIFaddbegin \DIFadd{Dissatisfaction in a foreign country is natural if the language native to that region is not known}\DIFaddend . \cite{ryan2001} also found that communication difficulty was one of the main reasons \DIFaddbegin \DIFadd{that }\DIFaddend Chinese customers would state for not visiting again. \DIFdelbegin \DIFdel{It seems like this is a problem that is not singular }\DIFdelend \DIFaddbegin \DIFadd{However, this issue is not exclusive }\DIFaddend to Japan.

    Our initial question was whether the environment of high-grade hospitality would affect both cultures equally. This study \DIFdelbegin \DIFdel{brought us closer to }\DIFdelend \DIFaddbegin \DIFadd{attempted to determine }\DIFaddend the answer. \DIFdelbegin \DIFdel{On the one hand, there is a possibility }\DIFdelend \DIFaddbegin \DIFadd{It is possible }\DIFaddend that Chinese customers \DIFdelbegin \DIFdel{did have }\DIFdelend \DIFaddbegin \DIFadd{had }\DIFaddend high-grade hospitality and \DIFdelbegin \DIFdel{did not get }\DIFdelend \DIFaddbegin \DIFadd{were }\DIFaddend equally satisfied with Westerners. In that case, it appears that the difference \DIFaddbegin \DIFadd{in perception }\DIFaddend stems from a psychological source\DIFdelbegin \DIFdel{. Expectation }\DIFdelend \DIFaddbegin \DIFadd{; expectation }\DIFaddend leads to satisfaction\DIFaddbegin \DIFadd{, }\DIFaddend and a lack of expectation results in lesser satisfaction. \DIFdelbegin \DIFdel{On the other hand, there }\DIFdelend \DIFaddbegin \DIFadd{There }\DIFaddend is also a possibility that Chinese customers are not receiving the highest grade of hospitality because of cultural friction between Japan and China.

    \DIFdelbegin \DIFdel{It }\DIFdelend \DIFaddbegin \DIFadd{The case representing the higher possibility }\DIFaddend is unclear from our results\DIFdelbegin \DIFdel{which of these could be the case. One thing is clear for hotel managers, however. Competing }\DIFdelend \DIFaddbegin \DIFadd{. However, competing }\DIFaddend in hospitality and service \DIFdelbegin \DIFdel{does include }\DIFdelend \DIFaddbegin \DIFadd{includes }\DIFaddend language services, especially in the international tourism industry. Better multilingual support can only improve \DIFdelbegin \DIFdel{that already high }\DIFdelend \DIFaddbegin \DIFadd{the hospitality }\DIFaddend standard in Japan. Considering that most of the tourists in Japan come from other countries in Asia, \DIFdelbegin \DIFdel{this is an endeavor that truly can bring benefits to their investment}\DIFdelend \DIFaddbegin \DIFadd{multilingual support is beneficial}\DIFaddend . Proposals for this endeavor include hiring Chinese speaking staff, preparing pamphlets in Chinese, or \DIFdelbegin \DIFdel{have }\DIFdelend \DIFaddbegin \DIFadd{having }\DIFaddend a translator application readily available with staff trained in interacting through an electronic translator.

  \subsection{Hard vs. soft satisfaction factors}\label{disc:hard_soft}

    As \DIFdelbegin \DIFdel{we }\DIFdelend stated in section \ref{theory_satisfaction}, previous research \DIFdelbegin \DIFdel{is focused mostly }\DIFdelend \DIFaddbegin \DIFadd{has mostly focused }\DIFaddend on the hotel's soft attributes and their influence on customer satisfaction. Examples of soft attributes include staff behavior, commodities, amenities, and appliances that can be improved within the hotel \DIFdelbegin %DIFDELCMD < \cite[e.g.][]{shanka2004,choi2001}%%%
\DIFdelend \DIFaddbegin \cite[e.g.,][]{shanka2004,choi2001}\DIFaddend . However, hard attributes are not usually analyzed in satisfaction studies. Examples of hard attributes include the hotel's location relative to public \DIFdelbegin \DIFdel{transport }\DIFdelend \DIFaddbegin \DIFadd{transportation }\DIFaddend and shops, language immersion of the country, noise pollution, or weather. Because \DIFdelbegin \DIFdel{our study left }\DIFdelend the satisfaction factors \DIFdelbegin \DIFdel{to be decided statistically }\DIFdelend \DIFaddbegin \DIFadd{were decided statistically in our study }\DIFaddend via customers’ online reviews, we can see the importance of \DIFdelbegin \DIFdel{those }\DIFdelend \DIFaddbegin \DIFadd{the }\DIFaddend hard or soft attributes in their priorities.

    Figure \ref{fig:hard_soft_zh} shows that\DIFaddbegin \DIFadd{, }\DIFaddend in regards to Chinese customer satisfaction, in general, 68 \% of the top 10 keywords are hard factors\DIFdelbegin \DIFdel{. In }\DIFdelend \DIFaddbegin \DIFadd{; in }\DIFaddend contrast, only 20 \% are soft factors. The rates are similar for most price ranges \DIFdelbegin \DIFdel{, excepting }\DIFdelend \DIFaddbegin \DIFadd{except }\DIFaddend the highest-priced hotels, where 35 \% of the keywords are undefined. However, the soft attributes are still similar at 18 \%. However, two of these managerial words are all concentrated at the top of the list (``\begin{CJK}{UTF8}{gbsn}不错\end{CJK} (not bad)\DIFdelbegin \DIFdel{'',}\DIFdelend \DIFaddbegin \DIFadd{,'' }\DIFaddend ``\begin{CJK}{UTF8}{gbsn}干净\end{CJK} (clean)''), \DIFdelbegin \DIFdel{plus }\DIFdelend \DIFaddbegin \DIFadd{and }\DIFaddend the adjective pairs \DIFdelbegin \DIFdel{relating }\DIFdelend \DIFaddbegin \DIFadd{related }\DIFaddend to soft attributes of ``\begin{CJK}{UTF8}{gbsn}不错\end{CJK} (not bad)'' \DIFdelbegin \DIFdel{which are }\DIFdelend \DIFaddbegin \DIFadd{are also }\DIFaddend at the top in most price ranges\DIFdelbegin \DIFdel{as well}\DIFdelend . Chinese tourists \DIFdelbegin \DIFdel{could }\DIFdelend \DIFaddbegin \DIFadd{may }\DIFaddend expect spaciousness and cleanliness when coming to Japan. \DIFdelbegin \DIFdel{That expectation could be caused by }\DIFdelend \DIFaddbegin \DIFadd{The expectation may be due to }\DIFaddend reputation, previous experiences, or cultural backgrounds. Some scholars argue that different cultures have different room size perceptions \cite[][]{Saulton2017}. Although the study subjects are German and South Korean, \DIFdelbegin \DIFdel{the study }\DIFdelend \DIFaddbegin \cite{Saulton2017} \DIFaddend presents the results as differences influenced by Asian and Western cultures. We argue that one country is not representative of others’ cultures, \DIFdelbegin \DIFdel{so }\DIFdelend \DIFaddbegin \DIFadd{and thus, }\DIFaddend there can be differences between South Korea and China in room size perception. However, \DIFdelbegin \DIFdel{an interesting point appears. It could be that }\DIFdelend a different room size perception \DIFdelbegin \DIFdel{affects }\DIFdelend \DIFaddbegin \DIFadd{may affect }\DIFaddend the satisfaction of Chinese tourists in contrast \DIFdelbegin \DIFdel{with Westerners . Westerners only start placing a priority on praising }\DIFdelend \DIFaddbegin \DIFadd{to Westerners because Westerners prioritize }\DIFaddend room size as the price of the hotel \DIFdelbegin \DIFdel{goes up}\DIFdelend \DIFaddbegin \DIFadd{room rises}\DIFaddend . We can compare these results with previous literature, where traveling Chinese tourists choose their destination based on several factors, including cleanliness, nature, architecture, and scenery \cite[][]{ryan2001}. These \DIFdelbegin \DIFdel{other few }\DIFdelend factors found in previous literature could be linked to the keyword ``\begin{CJK}{UTF8}{gbsn}环境\end{CJK} (environment or surroundings)'' as well. This keyword \DIFdelbegin \DIFdel{is present in }\DIFdelend \DIFaddbegin \DIFadd{was found for }\DIFaddend hotels priced at more than 20,000 yen per night. 

    In \DIFdelbegin \DIFdel{comparison}\DIFdelend \DIFaddbegin \DIFadd{contrast}\DIFaddend , English speakers are mostly satisfied with the hotel's soft attributes. Figure \ref{fig:hard_soft_en} shows that soft attributes are above 48 \% in all price ranges, the highest being 65 \% in the \DIFaddbegin \DIFadd{price range of }\DIFaddend 15,000 to 20,000 yen per night\DIFdelbegin \DIFdel{price range. This price range corresponds toaffordable business hotels}\DIFdelend \DIFaddbegin \DIFadd{, which corresponds to}\DIFaddend , for example\DIFaddbegin \DIFadd{, affordable business hotels}\DIFaddend . English-speaking customers also have soft attributes at the top of their list. The exception is the hard attribute that is the hotel's location, which is consistently around the middle of the top 10 lists for all price ranges. If one considers \DIFaddbegin \DIFadd{the satisfaction of }\DIFaddend both Chinese and Western tourists\DIFdelbegin \DIFdel{’ satisfaction}\DIFdelend , a hotel can improve \DIFaddbegin \DIFadd{its services }\DIFaddend to attract more customers in the future. \DIFdelbegin \DIFdel{If it was the other way around, and }\DIFdelend \DIFaddbegin \DIFadd{Otherwise, if }\DIFaddend the satisfaction was related more with hard attributes overall for 1020 both cultures, hotels \DIFdelbegin \DIFdel{would have to compete solely on their }\DIFdelend \DIFaddbegin \DIFadd{should be built considering the }\DIFaddend location.

    For both customer groups, the main reason for dissatisfaction \DIFdelbegin \DIFdel{is }\DIFdelend \DIFaddbegin \DIFadd{was }\DIFaddend pricing, which can be interpreted as a concern about value for money. However, \DIFdelbegin \DIFdel{it is interesting to note that while }\DIFdelend English-speaking customers \DIFdelbegin \DIFdel{complain }\DIFdelend \DIFaddbegin \DIFadd{complained }\DIFaddend about price with a lower rank in lower-priced hotels. In contrast, \DIFdelbegin \DIFdel{the }\DIFdelend Chinese customers consistently \DIFdelbegin \DIFdel{have }\DIFdelend \DIFaddbegin \DIFadd{had }\DIFaddend ``\begin{CJK}{UTF8}{gbsn}价格\end{CJK} (price)'' as \DIFdelbegin \DIFdel{a top }\DIFdelend \DIFaddbegin \DIFadd{the first }\DIFaddend or second-most concern across all price ranges. A \DIFdelbegin \DIFdel{paper studying }\DIFdelend \DIFaddbegin \DIFadd{study on }\DIFaddend Chinese tourists found that they had this concern \cite[][]{truong2009}. However, our results indicate that this \DIFdelbegin \DIFdel{is less of a }\DIFdelend \DIFaddbegin \DIFadd{concern is less related to the }\DIFaddend cultural attribute in Japanese hotels and \DIFdelbegin \DIFdel{has more to do }\DIFdelend \DIFaddbegin \DIFadd{more related }\DIFaddend with the pricing of hotels overall. \DIFdelbegin \DIFdel{The tourists }\DIFdelend \DIFaddbegin \DIFadd{Tourists }\DIFaddend coming to Japan could be \DIFdelbegin \DIFdel{both experienced travelers }\DIFdelend \DIFaddbegin \DIFadd{either experienced }\DIFaddend or first-time travelers. However, \DIFdelbegin \DIFdel{the fact is that }\DIFdelend their expectation of the price for hotels was lower than \DIFdelbegin \DIFdel{what they found }\DIFdelend \DIFaddbegin \DIFadd{the actual prices }\DIFaddend in Japan. In general, Japan is an expensive place to visit, \DIFaddbegin \DIFadd{thereby }\DIFaddend impacting this placement in the ranking. Space is scarce in Japan, and capsule hotels with cramped spaces of 2 x 1 meters cost around \DIFdelbegin \DIFdel{3,000 to 6,000 }\DIFdelend \DIFaddbegin \DIFadd{3000 to 6000 }\DIFaddend yen per night. Bigger business hotel rooms are relatively expensive, ranging from \DIFdelbegin \DIFdel{5,000 }\DIFdelend \DIFaddbegin \DIFadd{5000 }\DIFaddend to 12,000 yen per night. For comparison, hotels in the USA with a similar quality can \DIFdelbegin \DIFdel{be }\DIFdelend \DIFaddbegin \DIFadd{charge }\DIFaddend half the price.

    Around half of the dissatisfaction factors for both Chinese and Western customers are caused by issues that could be solved with improved management\DIFdelbegin \DIFdel{. The previous }\DIFdelend \DIFaddbegin \DIFadd{; this }\DIFaddend is true for all price ranges. \DIFdelbegin \DIFdel{Of course, the }\DIFdelend \DIFaddbegin \DIFadd{The }\DIFaddend improvements could be staff training (perhaps in language), hiring professional cleaning services for rooms with cigarette smoke smells, or improving the bedding\DIFdelbegin \DIFdel{. All of these options }\DIFdelend \DIFaddbegin \DIFadd{; however, all of these considerations }\DIFaddend can be costly. However, this paper provides a \DIFdelbegin \DIFdel{good guideline for which factors to consider first and which ones will be best suited to }\DIFdelend \DIFaddbegin \DIFadd{useful guideline about the factors to be prioritized and the factors that would be the most suitable for }\DIFaddend each customer group. Hotels can \DIFaddbegin \DIFadd{also }\DIFaddend use the price range categorization \DIFdelbegin \DIFdel{in order }\DIFdelend to choose the appropriate strategy\DIFdelbegin \DIFdel{as well}\DIFdelend . However, once the hotel's location and construction are set for Chinese customers, \DIFdelbegin \DIFdel{not much else can be done }\DIFdelend \DIFaddbegin \DIFadd{only a few changes can be made }\DIFaddend to satisfy them further. As mentioned \DIFdelbegin \DIFdel{before}\DIFdelend \DIFaddbegin \DIFadd{previously}\DIFaddend , Chinese language availability is another soft attribute that can be improved with staff and training investment.

    \DIFdelbegin \DIFdel{On the other hand, }\DIFdelend Western tourists are \DIFdelbegin \DIFdel{all around dissatisfied with mostly }\DIFdelend \DIFaddbegin \DIFadd{mainly dissatisfied with }\DIFaddend soft attributes. \DIFdelbegin \DIFdel{They show this by having a low }\DIFdelend \DIFaddbegin \DIFadd{This is revealed by a low satisfaction level }\DIFaddend of 35 \% in the highest price range where undefined factors are the majority and \DIFaddbegin \DIFadd{a maximum of }\DIFaddend 78 \% \DIFdelbegin \DIFdel{at most }\DIFdelend in the price range \DIFdelbegin \DIFdel{from }\DIFdelend \DIFaddbegin \DIFadd{of }\DIFaddend 30,000 to 50,000 yen per night in a hotel. \DIFdelbegin \DIFdel{The room for improvement for }\DIFdelend \DIFaddbegin \DIFadd{Improvement scope for }\DIFaddend Western tourists is more extensive than \DIFaddbegin \DIFadd{that for }\DIFaddend their Chinese counterparts. As such, it presents a \DIFdelbegin \DIFdel{bigger }\DIFdelend \DIFaddbegin \DIFadd{larger }\DIFaddend investment opportunity. As mentioned earlier in this paper, Westerners are known as ``long-haul'' customers, \DIFdelbegin \DIFdel{spending }\DIFdelend \DIFaddbegin \DIFadd{since they spend }\DIFaddend more than 45 \% of their budget on hotel lodging.
    \DIFdelbegin \DIFdel{On the other hand, their Asian counterparts only spend }\DIFdelend \DIFaddbegin \DIFadd{Asian tourists spend only }\DIFaddend 25 \% of their budget on hotels \cite[][]{choi2000}. With bigger returns on managerial improvements, \DIFdelbegin \DIFdel{it seems like we can }\DIFdelend \DIFaddbegin \DIFadd{we }\DIFaddend recommend investing in improving attributes that dissatisfy Western customers, such as cleanliness and removing tobacco smell. Making more hotel facilities tobacco-free and deodorizing the rooms can be a low-cost investment\DIFdelbegin \DIFdel{that }\DIFdelend \DIFaddbegin \DIFadd{, which }\DIFaddend could increase returns \DIFdelbegin \DIFdel{many timesover}\DIFdelend \DIFaddbegin \DIFadd{by several times}\DIFaddend .

    However, \DIFdelbegin \DIFdel{the opposite argument could also be made that Chinese customers provide a more significant numberof customers}\DIFdelend \DIFaddbegin \DIFadd{Chinese customers are more in number}\DIFaddend , even though they tend to spend less on lodging. Attracting a large number of Chinese customers can be a viable strategy for hotels. However, as mentioned before, they tend to focus more on hard attributes\DIFdelbegin \DIFdel{, leaving language barrier-breaking as }\DIFdelend \DIFaddbegin \DIFadd{; thus, breaking the language barrier breaking is }\DIFaddend one of the few strategies to accomplish this.

    The basic premise of this study is that different cultures lead to different expectations and satisfaction factors. This premise also plays a role in the differentiation between the \DIFdelbegin \DIFdel{preference }\DIFdelend \DIFaddbegin \DIFadd{preferences }\DIFaddend of hard or soft attributes.

    In \cite{donthu1998cultural}, subjects from 10 different countries were compared \DIFdelbegin \DIFdel{in }\DIFdelend \DIFaddbegin \DIFadd{with respect to }\DIFaddend their expectations of service quality and analyzed \DIFdelbegin \DIFdel{through the lens of }\DIFdelend \DIFaddbegin \DIFadd{based on }\DIFaddend Hofstede's typology of culture \cite[][]{hofstede1984culture}. \DIFdelbegin \DIFdel{That }\DIFdelend \DIFaddbegin \DIFadd{The previous }\DIFaddend study states that\DIFaddbegin \DIFadd{, }\DIFaddend although culture has no \DIFdelbegin \DIFdel{one }\DIFdelend specific index, five dimensions of culture can be used to analyze or categorize a country in comparison to others. These are \textit{power distance}, \textit{uncertainty avoidance}, \textit{\DIFdelbegin \DIFdel{individualism-collectivism}\DIFdelend \DIFaddbegin \DIFadd{individualism–-collectivism}\DIFaddend }, \textit{\DIFdelbegin \DIFdel{masculinity-femininity}\DIFdelend \DIFaddbegin \DIFadd{masculinity–femininity}\DIFaddend }, and \textit{\DIFdelbegin \DIFdel{long-term versus }\DIFdelend \DIFaddbegin \DIFadd{long-term-–}\DIFaddend short-term orientation}. In each of these dimensions, at least one element of service expectations was found to be significantly different for countries grouped under contrasting attributes (e.g., individualistic countries vs. collectivist countries, high uncertainty avoidance countries vs. low uncertainty avoidance countries). However, Hofstede's typology has received criticism from academics, particularly \DIFaddbegin \DIFadd{for }\DIFaddend the fifth dimension that Hofstede proposed, which was \DIFdelbegin \DIFdel{added afterward with the alternate }\DIFdelend \DIFaddbegin \DIFadd{later added with the alternative }\DIFaddend name \textit{Confucian dynamic}. Academics with a Chinese background criticized Hofstede for being misinformed on the philosophical aspects of Confucianism \DIFdelbegin \DIFdel{, }\DIFdelend as well as \DIFdelbegin \DIFdel{being }\DIFdelend \DIFaddbegin \DIFadd{considering }\DIFaddend a difficult dimension to measure \cite[][]{fang2003critique}. Other models, such as the GLOBE model, also \DIFdelbegin \DIFdel{take issue with }\DIFdelend \DIFaddbegin \DIFadd{consider }\DIFaddend some of Hofstede's dimensions and replace them with others, making a total of \DIFdelbegin \DIFdel{9 }\DIFdelend \DIFaddbegin \DIFadd{nine }\DIFaddend dimensions \cite[][]{house1999cultural}. The \textit{\DIFdelbegin \DIFdel{masculinity-femininity}\DIFdelend \DIFaddbegin \DIFadd{masculinity–-femininity}\DIFaddend } dimension, for example, is proposed to be instead \DIFaddbegin \DIFadd{of }\DIFaddend two dimensions: \textit{gender egalitarianism} and \textit{assertiveness}. This addition of dimensions avoids assuming that assertiveness is either masculine or feminine, which stems from outdated gender stereotypes. \DIFdelbegin \DIFdel{Gender stereotypes such as these }\DIFdelend \DIFaddbegin \DIFadd{Such gender stereotypes }\DIFaddend have also been the subject of critique \DIFdelbegin \DIFdel{for }\DIFdelend \DIFaddbegin \DIFadd{on }\DIFaddend Hofstede's model\cite[][]{jeknic2014gender}. \DIFdelbegin \DIFdel{Our study agrees }\DIFdelend \DIFaddbegin \DIFadd{We agree }\DIFaddend with these critiques and \DIFdelbegin \DIFdel{, therefore, will avoid considering these for }\DIFdelend \DIFaddbegin \DIFadd{thus avoid considering such stereotypes in }\DIFaddend our discussion.

    The backgrounds of collectivism in China and individualism in Western countries have been studied \DIFdelbegin \DIFdel{before }\DIFdelend \DIFaddbegin \DIFadd{previously }\DIFaddend \cite[][]{gao2017chinese}. These backgrounds \DIFdelbegin \DIFdel{and }\DIFdelend \DIFaddbegin \DIFadd{as well as }\DIFaddend the differences in these cultural dimensions could be the underlying cause for differences in expectations. Regardless of the cause, however, measures in the past have proven that \DIFdelbegin \DIFdel{these differences in expectations }\DIFdelend \DIFaddbegin \DIFadd{such differences }\DIFaddend exist \cite[][]{armstrong1997importance}. 

    For our purposes \DIFdelbegin \DIFdel{in }\DIFdelend \DIFaddbegin \DIFadd{of }\DIFaddend contrasting Western vs. Chinese satisfaction stemming from expectations, these dimensions could explain why Chinese customers are generally satisfied more often with hard factors while Westerners are satisfied or dissatisfied with soft factors. \DIFdelbegin \DIFdel{Perhaps the }\DIFdelend \DIFaddbegin \DIFadd{The }\DIFaddend cultural background of Chinese tourists emphasizes their surroundings and their place in nature and the environment. Chinese historical backgrounds of Confucianism, Taoism, and Buddhism permeate the thought processes of Chinese populations. However, scholars argue that the changes in generations and their economic and recent history \DIFdelbegin \DIFdel{gives }\DIFdelend \DIFaddbegin \DIFadd{attaches }\DIFaddend less importance to these concepts in their lives \cite[][]{gao2017chinese}. 

    Nevertheless, \DIFdelbegin \DIFdel{one could argue that }\DIFdelend a Chinese cultural attribute emphasizes \DIFaddbegin \DIFadd{that }\DIFaddend the environment and the \DIFdelbegin \DIFdel{place one is in towards }\DIFdelend \DIFaddbegin \DIFadd{location  affect }\DIFaddend satisfaction, rather than the \DIFdelbegin \DIFdel{way one is treated}\DIFdelend \DIFaddbegin \DIFadd{treatment}\DIFaddend . According to previous research, Chinese tourists are \DIFdelbegin \DIFdel{collectivist, while }\DIFdelend \DIFaddbegin \DIFadd{collectivists, whereas }\DIFaddend Westerners are individualists \cite[][]{kim2000}. A more anthropocentric and individualistic Western culture could \DIFdelbegin \DIFdel{result in }\DIFdelend \DIFaddbegin \DIFadd{correlate }\DIFaddend more of their expectations and priorities \DIFdelbegin \DIFdel{be related to how one is treated }\DIFdelend \DIFaddbegin \DIFadd{to the treatment }\DIFaddend in social circumstances, rather than the environment\DIFdelbegin \DIFdel{one is in}\DIFdelend . According to \cite{donthu1998cultural}, highly individualistic customers\DIFaddbegin \DIFadd{, in contrast to collectivistic customers, }\DIFaddend have a higher expectation of empathy and assurance from the provider\DIFdelbegin \DIFdel{than collectivist customers. Empathy and assurance from the provider }\DIFdelend \DIFaddbegin \DIFadd{, which }\DIFaddend are aspects of service, a soft attribute of a hotel.

    Among other dimensions in both models, we can consider uncertainty avoidance. \DIFdelbegin \DIFdel{High uncertainty avoidance customers would }\DIFdelend \DIFaddbegin \DIFadd{Customers of high uncertainty avoidance }\DIFaddend carefully plan their travel and \DIFdelbegin \DIFdel{therefore }\DIFdelend \DIFaddbegin \DIFadd{thus }\DIFaddend have higher expectations towards service. \DIFdelbegin \DIFdel{On the other hand, }\DIFdelend \DIFaddbegin \DIFadd{In contrast, customers of }\DIFaddend lower uncertainty avoidance \DIFdelbegin \DIFdel{customers would have certain room for risk }\DIFdelend \DIFaddbegin \DIFadd{do not take risks }\DIFaddend in their decisions \DIFdelbegin \DIFdel{, and therefore }\DIFdelend \DIFaddbegin \DIFadd{and thus }\DIFaddend face less disappointment with different expectations. However, according to \cite{xiumei2011cultural}, the difference between China and the USA in uncertainty avoidance is not \DIFdelbegin \DIFdel{so }\DIFdelend clear when measuring with the Hofstede typology and the GLOBE typology. While the USA is not representative of Western society, \DIFdelbegin \DIFdel{this dimension might not be the one causing the difference in hard-soft }\DIFdelend \DIFaddbegin \DIFadd{uncertainty avoidance may not cause the difference in hard–-soft }\DIFaddend attribute satisfaction between Chinese and Western cultures. \DIFdelbegin \DIFdel{Another }\DIFdelend \DIFaddbegin \DIFadd{Differences in another }\DIFaddend factor, power distance, \DIFdelbegin \DIFdel{was also different when measured }\DIFdelend \DIFaddbegin \DIFadd{were also noted when measuring }\DIFaddend by Hofstede's method compared to the GLOBAL method\DIFdelbegin \DIFdel{, so we decided against making this }\DIFdelend \DIFaddbegin \DIFadd{; therefore, power distance was not considered for }\DIFaddend comparison.

  \subsection{Satisfaction across different price ranges}\label{disc:price}

    In previous sections of this paper, we \DIFdelbegin \DIFdel{have }\DIFdelend mentioned the differences reflected in hotel price ranges. \DIFdelbegin \DIFdel{Nevertheless, it is interesting to discuss this further. }\DIFdelend The most visible change \DIFdelbegin \DIFdel{in }\DIFdelend \DIFaddbegin \DIFadd{of }\DIFaddend satisfaction factors across differently priced hotels is the change in voice \DIFdelbegin \DIFdel{to describe }\DIFdelend \DIFaddbegin \DIFadd{when describing }\DIFaddend their satisfaction with the same topics. We \DIFdelbegin \DIFdel{can know }\DIFdelend \DIFaddbegin \DIFadd{noticed }\DIFaddend this by observing the adjective\DIFdelbegin \DIFdel{+ noun }\DIFdelend \DIFaddbegin \DIFadd{–-noun }\DIFaddend pairs and finding pairs with different adjectives for the same nouns. For example, in English, words describing nouns such as ``location'' or ``hotel'' are ``good'' or ``nice'' in lower-priced hotels. In contrast, the adjectives that pair with the same nouns for \DIFdelbegin \DIFdel{more highly-priced }\DIFdelend \DIFaddbegin \DIFadd{higher- priced }\DIFaddend hotels are ``wonderful'' and ``excellent\DIFdelbegin \DIFdel{''.}\DIFdelend \DIFaddbegin \DIFadd{.'' }\DIFaddend In Chinese, the change \DIFdelbegin \DIFdel{goes }\DIFdelend \DIFaddbegin \DIFadd{ranges }\DIFaddend from ``\begin{CJK}{UTF8}{gbsn}不错\end{CJK} (not bad)'' to ``\begin{CJK}{UTF8}{gbsn}棒\end{CJK} (great)'' or ``\begin{CJK}{UTF8}{gbsn}赞\end{CJK} (awesome)\DIFdelbegin \DIFdel{''.}\DIFdelend \DIFaddbegin \DIFadd{.'' }\DIFaddend We can infer that the level of satisfaction is high and influences how customers write their reviews. \DIFdelbegin \DIFdel{However, when we look at }\DIFdelend \DIFaddbegin \DIFadd{Regarding }\DIFaddend the negative keywords, \DIFdelbegin \DIFdel{the change is }\DIFdelend \DIFaddbegin \DIFadd{however, the change ranges }\DIFaddend from ``annoying'' or ``worst'' \DIFdelbegin \DIFdel{, }\DIFdelend to ``disappointing\DIFdelbegin \DIFdel{''.Here we can see }\DIFdelend \DIFaddbegin \DIFadd{.'' Here, we can determine }\DIFaddend how expectations influence satisfaction and dissatisfaction in different ways. 

    In this paper, we follow the definition of satisfaction by \cite{hunt1975}, where meeting or exceeding expectations produces satisfaction. Therefore, the \DIFdelbegin \DIFdel{lack }\DIFdelend \DIFaddbegin \DIFadd{failure }\DIFaddend of meeting expectations would cause dissatisfaction. In the \DIFdelbegin \DIFdel{cases above}\DIFdelend \DIFaddbegin \DIFadd{aforementioned cases}\DIFaddend , we can infer that a customer that pays more for a \DIFdelbegin \DIFdel{higher class of }\DIFdelend \DIFaddbegin \DIFadd{higher-class }\DIFaddend experience has higher expectations. This is true in dissatisfaction, where their expectation is higher in a more expensive hotel. As such, any lack of cleanliness can lead to disappointment\DIFdelbegin \DIFdel{or outrage}\DIFdelend . In the case of English-speaking customers in the 30,\DIFdelbegin \DIFdel{000 to }\DIFdelend \DIFaddbegin \DIFadd{000-–}\DIFaddend 50,000 yen per night price range, cigarette smell is particularly disappointing. However, we consistently see customers with high expectations for high-class hotels \DIFdelbegin \DIFdel{react }\DIFdelend \DIFaddbegin \DIFadd{reacting }\DIFaddend even more positively when satisfied. In the positive case, expectations appear to be exceeded in most cases, judging \DIFaddbegin \DIFadd{from }\DIFaddend their reactions. We argue that these are two different kinds of \DIFdelbegin \DIFdel{interactions with expectations. We can observe logical expectations. Customers set a standard in their mind }\DIFdelend \DIFaddbegin \DIFadd{expectations: logical and emotional. In the first case, customers are determined }\DIFaddend that the service must not fall below \DIFdelbegin \DIFdel{or be disappointed — }\DIFdelend \DIFaddbegin \DIFadd{a specific standard; }\DIFaddend for example, \DIFdelbegin \DIFdel{a customer being disappointed with dirty }\DIFdelend \DIFaddbegin \DIFadd{they can be disappointed with unhygienic }\DIFaddend rooms or cigarette smell.

    In \DIFdelbegin \DIFdel{contrast, we can observe emotional expectations, where }\DIFdelend \DIFaddbegin \DIFadd{the second case, in contrast, }\DIFaddend customers have a vague idea of having a positive experience \DIFdelbegin \DIFdel{. However, they }\DIFdelend \DIFaddbegin \DIFadd{but }\DIFaddend do not measure it against any standard. For example, \DIFdelbegin \DIFdel{having }\DIFdelend \DIFaddbegin \DIFadd{they expect }\DIFaddend a pleasant customer service experience or \DIFdelbegin \DIFdel{being treated hospitably }\DIFdelend \DIFaddbegin \DIFadd{a hospitable treatment }\DIFaddend by the staff at a high-class hotel. Regardless of their knowledge \DIFdelbegin \DIFdel{beforehand of the service to be provided}\DIFdelend \DIFaddbegin \DIFadd{in advance}\DIFaddend , positive emotions \DIFdelbegin \DIFdel{give }\DIFdelend \DIFaddbegin \DIFadd{offer }\DIFaddend them a perception of exceeded expectations and high satisfaction. \DIFdelbegin \DIFdel{This is where }\DIFdelend \DIFaddbegin \DIFadd{Thus, }\DIFaddend hospitality and service \DIFdelbegin \DIFdel{come into play and enhances }\DIFdelend \DIFaddbegin \DIFadd{enhance }\DIFaddend the experience of the customers. 

    There are interesting differences between Chinese and English-speaking tourists in their \DIFdelbegin \DIFdel{change in }\DIFdelend satisfaction factors to differently priced hotels. For example, \DIFdelbegin \DIFdel{we can observe that the }\DIFdelend Chinese tourists have ``\begin{CJK}{UTF8}{gbsn}购物\end{CJK} (shopping)'' as a top keyword in all the price ranges. In contrast, English-speaking tourists \DIFdelbegin \DIFdel{only mention it }\DIFdelend \DIFaddbegin \DIFadd{mention it only }\DIFaddend as a top keyword in the 20,000\DIFdelbegin \DIFdel{to 30}\DIFdelend \DIFaddbegin \DIFadd{–-30}\DIFaddend ,000 yen price range. It is \DIFdelbegin \DIFdel{common knowledge }\DIFdelend \DIFaddbegin \DIFadd{widely known }\DIFaddend in Japan that \DIFdelbegin \DIFdel{Chinese tourists coming to Japan with the express intention of shoppingare common}\DIFdelend \DIFaddbegin \DIFadd{many Chinese tourists visit Japan for shopping}\DIFaddend . \cite{tsujimoto2017purchasing} analyzed the souvenir purchasing behavior of Chinese tourists in Japan. The study shows that common products besides food and drink are: electronics, cameras, cosmetics, \DIFaddbegin \DIFadd{and }\DIFaddend medicine, among other more traditional \textit{souvenir} items, such as objects \DIFaddbegin \DIFadd{that are }\DIFaddend representative of the culture or places \DIFaddbegin \DIFadd{that }\DIFaddend they visit \cite{japan2014consumption}. \DIFdelbegin \DIFdel{There is an understanding that touristschoose }\DIFdelend \DIFaddbegin \DIFadd{Furthermore, tourists’ choice }\DIFaddend to shop in Japan \DIFdelbegin \DIFdel{has more to do }\DIFdelend \DIFaddbegin \DIFadd{is more related }\DIFaddend with the quality of the items \DIFaddbegin \DIFadd{rather }\DIFaddend than their relation to the tourist attractions. Our results \DIFdelbegin \DIFdel{suggest }\DIFdelend \DIFaddbegin \DIFadd{suggested }\DIFaddend that Western tourists \DIFdelbegin \DIFdel{are }\DIFdelend \DIFaddbegin \DIFadd{were }\DIFaddend engaging more in tourist attractions \DIFdelbegin \DIFdel{in comparison with }\DIFdelend \DIFaddbegin \DIFadd{rather than }\DIFaddend shopping activities. Another interesting difference is that English-speaking tourists start using negative keywords about the hotel's price only \DIFdelbegin \DIFdel{after }\DIFdelend \DIFaddbegin \DIFadd{if }\DIFaddend it concerns hotels of 15 \DIFdelbegin \DIFdel{,}\DIFdelend 000 yen or more\DIFdelbegin \DIFdel{, and it rises in its ranking }\DIFdelend \DIFaddbegin \DIFadd{; thereafter, }\DIFaddend the more expensive the hotel\DIFdelbegin \DIFdel{is}\DIFdelend \DIFaddbegin \DIFadd{, the higher the ranking}\DIFaddend . In contrast, \DIFdelbegin \DIFdel{Chinese customers have this keyword as their }\DIFdelend \DIFaddbegin \DIFadd{for Chinese customers, this keyword is the }\DIFaddend top keyword across all price ranges. Previous research suggests that value for money is a key concern for Chinese and Asian tourists \cite[][]{choi2000,choi2001,truong2009}, \DIFdelbegin \DIFdel{while }\DIFdelend \DIFaddbegin \DIFadd{whereas }\DIFaddend Western customers are more concerned \DIFdelbegin \DIFdel{with }\DIFdelend \DIFaddbegin \DIFadd{about }\DIFaddend hospitality \cite[][]{kozak2002}.  

    While some aspects of satisfaction and dissatisfaction change depending on the hotel's price range, some other factors \DIFdelbegin \DIFdel{stay mostly }\DIFdelend \DIFaddbegin \DIFadd{remain }\DIFaddend constant for each culture's customers. For example, appreciation for staff from English-speaking tourists is ranked close to the top satisfaction factor in all the price ranges. Satisfaction for cleanliness by both cultures constantly \DIFdelbegin \DIFdel{stays }\DIFdelend \DIFaddbegin \DIFadd{remains }\DIFaddend part of the top 10 keywords, except for the most expensive one, where other keywords \DIFdelbegin \DIFdel{take that place }\DIFdelend \DIFaddbegin \DIFadd{replace keywords related to satisfaction or cleanliness }\DIFaddend in the ranking\DIFdelbegin \DIFdel{. However, it is }\DIFdelend \DIFaddbegin \DIFadd{; however, they remain }\DIFaddend still high on the list. Chinese tourists have a high ranking for the word ``\begin{CJK}{UTF8}{gbsn}早餐\end{CJK} (breakfast)'' across all price ranges as well. As discussed in section \ref{disc:location}, \DIFdelbegin \DIFdel{transport }\DIFdelend \DIFaddbegin \DIFadd{transportation }\DIFaddend and location are also important for hotels of all classes and prices. While the ranking of attributes might differ between price ranges, hard and soft attribute proportions also appear to be constant within \DIFdelbegin \DIFdel{at most }\DIFdelend a 13 \% margin of error per attribute\DIFdelbegin \DIFdel{, often being lower}\DIFdelend . This suggests that\DIFdelbegin \DIFdel{culturally }\DIFdelend \DIFaddbegin \DIFadd{, from a cultural aspect, }\DIFaddend the customers have a particular bias to consider some attributes more than others.

  \subsection{Implications for hotel managers}\label{disc:implications}

    Our study \DIFdelbegin \DIFdel{presents }\DIFdelend \DIFaddbegin \DIFadd{reached }\DIFaddend two important conclusions: one about hospitality and cultural differences \DIFdelbegin \DIFdel{, }\DIFdelend and another about managerial decisions towards two different populations. \DIFdelbegin \DIFdel{As a whole}\DIFdelend \DIFaddbegin \DIFadd{Overall}\DIFaddend , Chinese tourists \DIFdelbegin \DIFdel{are not showing the most satisfaction towards }\DIFdelend \DIFaddbegin \DIFadd{did not attach much importance to }\DIFaddend hospitality and service factors. Instead, they \DIFdelbegin \DIFdel{focus }\DIFdelend \DIFaddbegin \DIFadd{focused }\DIFaddend on the hard attributes of a hotel. \DIFdelbegin \DIFdel{Either they do not get as much satisfaction from hospitality }\DIFdelend \DIFaddbegin \DIFadd{In particular, they were not satisfied with hospitality as much }\DIFaddend as Western tourists \DIFdelbegin \DIFdel{or feel }\DIFdelend \DIFaddbegin \DIFadd{were; otherwise, they felt }\DIFaddend that basic language and communication needs \DIFdelbegin \DIFdel{are not being met, so they receive a lesser impression. On the other hand, Western tourists are elated }\DIFdelend \DIFaddbegin \DIFadd{were not met, thereby they were not much satisfied. Western tourists were highly satisfied }\DIFaddend with Japanese hospitality \DIFdelbegin \DIFdel{, preferring }\DIFdelend \DIFaddbegin \DIFadd{and preferred }\DIFaddend soft attributes to \DIFdelbegin \DIFdel{hard-set }\DIFdelend \DIFaddbegin \DIFadd{hard }\DIFaddend ones. The other conclusion is that managerial decisions \DIFdelbegin \DIFdel{will }\DIFdelend \DIFaddbegin \DIFadd{could }\DIFaddend mostly benefit Western tourists, except that Chinese language improvements and breakfast inclusion \DIFdelbegin \DIFdel{can satisfy more Chinese customers . Japan is recently seeing }\DIFdelend \DIFaddbegin \DIFadd{could satisfy Chinese customers to a greater extent. Recently, Japan has been facing }\DIFaddend an increase in Chinese students as well as \DIFdelbegin \DIFdel{Western students of }\DIFdelend \DIFaddbegin \DIFadd{students of Western }\DIFaddend universities. Hiring students as part-time workers could increase the language services of a hotel.

    To satisfy both customer types, hotel managers need to invest in cleanliness, deodorizing, and making hotel rooms tobacco-free. It could also be recommended to invest in breakfast inclusion and multilingual services and staff preparedness to deal with Chinese and English speakers. Western tourists were also observed to have high comfort standards, which \DIFdelbegin \DIFdel{can be improved upon managerially }\DIFdelend \DIFaddbegin \DIFadd{could be managerially improved }\DIFaddend for better reviews. Perhaps it could be suggested to perform surveys of the bedding that is most comfortable for Western tourists. However, not all hotels can invest in all of these factors simultaneously. Our results suggest that satisfying cleanliness needs \DIFdelbegin \DIFdel{can }\DIFdelend \DIFaddbegin \DIFadd{could }\DIFaddend satisfy both customer types. A low-cost investment could be to make the facilities tobacco-free. Our results are also divided by price ranges, \DIFdelbegin \DIFdel{so }\DIFdelend \DIFaddbegin \DIFadd{thereby }\DIFaddend a hotel manager \DIFdelbegin \DIFdel{can }\DIFdelend \DIFaddbegin \DIFadd{could }\DIFaddend consider which analysis suits their hotel the most.

    \DIFdelbegin \DIFdel{While }\DIFdelend \DIFaddbegin \DIFadd{Albeit }\DIFaddend not manageable after a hotel \DIFdelbegin \DIFdel{has finished its construction}\DIFdelend \DIFaddbegin \DIFadd{is completely constructed}\DIFaddend , hard attributes \DIFdelbegin \DIFdel{are essential to consider for }\DIFdelend \DIFaddbegin \DIFadd{should be considered by }\DIFaddend managers. As \DIFdelbegin \DIFdel{previously stated , transport }\DIFdelend \DIFaddbegin \DIFadd{stated before, transportation }\DIFaddend systems and nearby shops are points of interest for both Chinese and Western tourists. Japanese hotel managers \DIFdelbegin \DIFdel{have to }\DIFdelend \DIFaddbegin \DIFadd{should }\DIFaddend consider the location and surroundings before the hotel is constructed. A suggestion could be to purchase land and start the construction \DIFdelbegin \DIFdel{after public plans to make }\DIFdelend \DIFaddbegin \DIFadd{at places where public plans for }\DIFaddend new subway lines \DIFdelbegin \DIFdel{are }\DIFdelend \DIFaddbegin \DIFadd{have already been }\DIFaddend made. 

    \DIFdelbegin \DIFdel{It is left to the managers to }\DIFdelend \DIFaddbegin \DIFadd{The managers must }\DIFaddend consider their business model for \DIFaddbegin \DIFadd{implementing }\DIFaddend the next strategy. One option could be attracting more Chinese customers \DIFdelbegin \DIFdel{in number }\DIFdelend with their observed low budgeting. Another could be attracting more \DIFdelbegin \DIFdel{high budget Western customersis on par with their business model}\DIFdelend \DIFaddbegin \DIFadd{big-budget Western customers}\DIFaddend . For example, investing more in cleanliness could improve Western customers looking for high-quality lodging satisfaction, even \DIFdelbegin \DIFdel{though the }\DIFdelend \DIFaddbegin \DIFadd{for an increased }\DIFaddend price per night\DIFdelbegin \DIFdel{would increase}\DIFdelend . On the other hand, hotels might be \DIFdelbegin \DIFdel{considered }\DIFdelend \DIFaddbegin \DIFadd{deemed }\DIFaddend costly by Chinese customers wherever such an investment is made.

\section{Limitations and Future Work}\label{limitations}

  \DIFdelbegin \DIFdel{This paper is not without its limitations. We }\DIFdelend \DIFaddbegin \DIFadd{In this study, we }\DIFaddend analyzed satisfaction and dissatisfaction keywords based on whether they appeared on satisfied reviews or dissatisfied ones. Following that, we attempted to understand the context in which these words were used by using a dependency parser and observing the related nouns. However, \DIFdelbegin \DIFdel{the }\DIFdelend \DIFaddbegin \DIFadd{our }\DIFaddend study is limited \DIFdelbegin \DIFdel{in that it only analyzes }\DIFdelend \DIFaddbegin \DIFadd{because it analyzed solely }\DIFaddend the words directly related to each keyword and \DIFdelbegin \DIFdel{does }\DIFdelend \DIFaddbegin \DIFadd{did }\DIFaddend not follow the upstream or downstream path down \DIFaddbegin \DIFadd{to }\DIFaddend further connections. This means that if the words \DIFdelbegin \DIFdel{are }\DIFdelend \DIFaddbegin \DIFadd{were }\DIFaddend used in combination with other keywords, we did not trace the effects of multiple contradicting statements. For example, in the sentence ``The room is good, but the food is lacking\DIFdelbegin \DIFdel{'',we would extract ``good food}\DIFdelend \DIFaddbegin \DIFadd{,'' we extracted ``good room}\DIFaddend '' and ``lacking food'' \DIFdelbegin \DIFdel{, but do }\DIFdelend \DIFaddbegin \DIFadd{but did }\DIFaddend not consider the fact that both occurred in the same sentence.

  This study analyzed the differences in customers' expectations at different levels of hospitality and service factors by dividing our data into price ranges. However, in the same price range, for example, the highest one, we can find both a \DIFdelbegin \DIFdel{western-style }\DIFdelend \DIFaddbegin \DIFadd{Western-style }\DIFaddend five-star resort and a high-end \DIFdelbegin \DIFdel{Japanese-style }\DIFdelend \DIFaddbegin \DIFadd{Japanese style }\DIFaddend \textit{ryokan}. Services offered in these hotels are \DIFdelbegin \DIFdel{very }\DIFdelend \DIFaddbegin \DIFadd{of }\DIFaddend high quality, \DIFdelbegin \DIFdel{although }\DIFdelend \DIFaddbegin \DIFadd{albeit }\DIFaddend very different. However, most of our database \DIFdelbegin \DIFdel{is }\DIFdelend \DIFaddbegin \DIFadd{was }\DIFaddend focused on the middle range priced hotels, \DIFdelbegin \DIFdel{which is }\DIFdelend \DIFaddbegin \DIFadd{the services of which are }\DIFaddend comparably less varied\DIFdelbegin \DIFdel{in service}\DIFdelend . However, there is still a \DIFdelbegin \DIFdel{divide between western }\DIFdelend \DIFaddbegin \DIFadd{division between Western }\DIFaddend and Japanese style hotels\DIFaddbegin \DIFadd{, which constitutes another limitation of our study}\DIFaddend .

  An essential aspect of this study is that we \DIFdelbegin \DIFdel{focus }\DIFdelend \DIFaddbegin \DIFadd{focused }\DIFaddend on the satisfaction and dissatisfaction towards \DIFdelbegin \DIFdel{expectations of }\DIFdelend \DIFaddbegin \DIFadd{the expectations of the }\DIFaddend individual aspects of the hotels. This \DIFdelbegin \DIFdel{gives }\DIFdelend \DIFaddbegin \DIFadd{gave }\DIFaddend us insight into \DIFdelbegin \DIFdel{which factors can focus on }\DIFdelend \DIFaddbegin \DIFadd{the factors that can be considered }\DIFaddend by hotel managers\DIFdelbegin \DIFdel{in applying this knowledge}\DIFdelend . However, \DIFdelbegin \DIFdel{our study was limited in that }\DIFdelend the overall satisfaction of each customer was not measured. This \DIFdelbegin \DIFdel{could be done }\DIFdelend \DIFaddbegin \DIFadd{measurement can be performed }\DIFaddend by rating the volume of text used to describe satisfaction \DIFdelbegin \DIFdel{factors against the text volume used for dissatisfaction }\DIFdelend \DIFaddbegin \DIFadd{and dissatisfaction factors}\DIFaddend . However, this \DIFdelbegin \DIFdel{imposes }\DIFdelend \DIFaddbegin \DIFadd{would impose }\DIFaddend a few difficulties\DIFdelbegin \DIFdel{that }\DIFdelend \DIFaddbegin \DIFadd{, which }\DIFaddend are out of the scope of this study. 

  Another limitation is that a large portion of \DIFdelbegin \DIFdel{the }\DIFdelend Asian tourists coming to Japan \DIFdelbegin \DIFdel{is }\DIFdelend \DIFaddbegin \DIFadd{are }\DIFaddend Taiwanese and Korean. We could not analyze these populations because our team members \DIFdelbegin \DIFdel{do }\DIFdelend \DIFaddbegin \DIFadd{did }\DIFaddend not know those languages. \DIFdelbegin \DIFdel{Aside from that}\DIFdelend \DIFaddbegin \DIFadd{Moreover}\DIFaddend , further typology analysis could not be made because of the nature of the data collected (for example, Chinese men and women of different ages or their Westerner counterparts).

  In future work, we plan to investigate \DIFdelbegin \DIFdel{further into this topic}\DIFdelend \DIFaddbegin \DIFadd{these topics further}\DIFaddend . We plan to extend our data to research \DIFdelbegin \DIFdel{for different trends for different }\DIFdelend \DIFaddbegin \DIFadd{of different trends and }\DIFaddend regions of Japan\DIFdelbegin \DIFdel{and }\DIFdelend \DIFaddbegin \DIFadd{, }\DIFaddend different kinds of hotels\DIFdelbegin \DIFdel{and between }\DIFdelend \DIFaddbegin \DIFadd{, and }\DIFaddend customers traveling alone or in groups, \DIFaddbegin \DIFadd{whether }\DIFaddend for fun or for work. Another point of interest in this study's future \DIFaddbegin \DIFadd{work }\DIFaddend is to use word clusters with similar meanings instead of single words. 

\section{Conclusion}\label{conclusion}

  In this study, our objective was to analyze the differences in satisfaction and dissatisfaction between Chinese and English-speaking customers of Japanese hotels, particularly in the context of Japanese hospitality, \textit{omotenashi}. To answer our research questions \ref{rsq:hospitality} and \ref{rsq:hospitality_both}, we extracted keywords from their online reviews uploaded to the portal sites \textit{Ctrip} and \textit{TripAdvisor} using Shannon's entropy calculations. We used these keywords for sentiment classification via an SVC. We then used dependency parsing and part of speech tagging to extract commonly found pairs of adjectives and nouns \DIFdelbegin \DIFdel{, }\DIFdelend as well as single words. We divided \DIFdelbegin \DIFdel{this }\DIFdelend \DIFaddbegin \DIFadd{these }\DIFaddend data by sentiment and hotel price range \DIFdelbegin \DIFdel{, }\DIFdelend \DIFaddbegin \DIFadd{by }\DIFaddend considering the most expensive room for one night. 

  In the context of Japanese hospitality, we found that Western tourists \DIFdelbegin \DIFdel{had the most satisfaction }\DIFdelend \DIFaddbegin \DIFadd{were most satisfied }\DIFaddend with staff behavior, cleanliness, and other attributes \DIFdelbegin \DIFdel{relating }\DIFdelend \DIFaddbegin \DIFadd{related }\DIFaddend to the hotel's services and hospitality. However, we found that Chinese customers had \DIFdelbegin \DIFdel{other concerns }\DIFdelend \DIFaddbegin \DIFadd{concerns other }\DIFaddend than hospitality when studying their satisfaction\DIFdelbegin \DIFdel{, }\DIFdelend \DIFaddbegin \DIFadd{; they were }\DIFaddend more inclined to praise the room, location, \DIFdelbegin \DIFdel{or }\DIFdelend \DIFaddbegin \DIFadd{and }\DIFaddend hotel's convenience. We found that \DIFdelbegin \DIFdel{both cultures have a different reaction to this hospitality environment . Both cultures have a different way of reacting to different prices. From this}\DIFdelend \DIFaddbegin \DIFadd{the two cultures had different reactions to the hospitality environment and the prices. Thus}\DIFaddend , we discussed two possible theories on \DIFaddbegin \DIFadd{the reasons }\DIFaddend why Chinese tourists \DIFdelbegin \DIFdel{respond differently to Westerners in this }\DIFdelend \DIFaddbegin \DIFadd{responded differently from Westerners in the }\DIFaddend environment of \textit{omotenashi}. One theory is that\DIFdelbegin \DIFdel{while they are being }\DIFdelend \DIFaddbegin \DIFadd{, although they were }\DIFaddend treated well and \DIFdelbegin \DIFdel{react }\DIFdelend \DIFaddbegin \DIFadd{thus reacted }\DIFaddend positively, the environment \DIFdelbegin \DIFdel{is }\DIFdelend \DIFaddbegin \DIFadd{was }\DIFaddend not compatible with them because of language or culture barriers, which \DIFdelbegin \DIFdel{lessens }\DIFdelend \DIFaddbegin \DIFadd{deteriorated }\DIFaddend their experience. The second possible theory is that they \DIFdelbegin \DIFdel{react differently to hospitality since they do }\DIFdelend \DIFaddbegin \DIFadd{reacted to hospitality differently since they did }\DIFaddend not have the same expectations\DIFdelbegin \DIFdel{to be satisfied in the same way. We theorize }\DIFdelend \DIFaddbegin \DIFadd{. We theorized }\DIFaddend that a lack of expectations could result in lessened satisfaction, even if the same service \DIFdelbegin \DIFdel{is presented}\DIFdelend \DIFaddbegin \DIFadd{was present}\DIFaddend . On the other hand, even when they \DIFdelbegin \DIFdel{hold }\DIFdelend \DIFaddbegin \DIFadd{held }\DIFaddend high expectations in a \DIFdelbegin \DIFdel{highly-priced hotel, Western tourists show that Japanese hospitality exceeds their }\DIFdelend \DIFaddbegin \DIFadd{high-priced hotel, Japanese hospitality exceeded Western tourists’ }\DIFaddend expectations, judging by their vocabulary for expressing their satisfaction. We \DIFdelbegin \DIFdel{consider }\DIFdelend \DIFaddbegin \DIFadd{considered }\DIFaddend that Western tourists \DIFdelbegin \DIFdel{are }\DIFdelend \DIFaddbegin \DIFadd{were }\DIFaddend more reactive to hospitality and service factors than their Chinese counterparts.

  Lastly, we measured the satisfaction and dissatisfaction factors, \DIFdelbegin \DIFdel{referring to }\DIFdelend \DIFaddbegin \DIFadd{that is, }\DIFaddend a hotel's hard and soft attributes. Soft attributes can be changed via management and staff by an improvement in services. \DIFdelbegin \DIFdel{On the other hand}\DIFdelend \DIFaddbegin \DIFadd{In contrast}\DIFaddend , hard attributes are physical and impractical elements to change, such as the size of a room that has already been constructed, the location of a hotel, closeness to convenient spots, or elements out of the control of the hotel managers. We found that\DIFaddbegin \DIFadd{, }\DIFaddend for satisfaction, Western tourists \DIFdelbegin \DIFdel{favor soft attributes . In contrast , Chinese touristsare }\DIFdelend \DIFaddbegin \DIFadd{favored soft attributes in contrast to Chinese tourists, who were }\DIFaddend more interested in the hard attributes of hotels \DIFdelbegin \DIFdel{, consistently across price ranges}\DIFdelend \DIFaddbegin \DIFadd{across all the price ranges consistently}\DIFaddend . For dissatisfaction, Western tourists \DIFdelbegin \DIFdel{are }\DIFdelend \DIFaddbegin \DIFadd{were }\DIFaddend also highly inclined to criticize soft attributes, such as cleanliness or cigarette smell in rooms. In contrast, Chinese tourists' dissatisfaction \DIFdelbegin \DIFdel{comes evenly }\DIFdelend \DIFaddbegin \DIFadd{derived }\DIFaddend from both hard and soft attributes \DIFaddbegin \DIFadd{evenly}\DIFaddend .

  One possible approach for hotel managers is to improve the satisfaction levels of Chinese tourists, who dedicate \DIFdelbegin \DIFdel{less }\DIFdelend \DIFaddbegin \DIFadd{lower }\DIFaddend percentage of their budget to hotels but are more abundant in number. They are less satisfied with soft attributes but have an identifiable method \DIFdelbegin \DIFdel{of }\DIFdelend \DIFaddbegin \DIFadd{for }\DIFaddend improving satisfaction by lessening language barriers and providing a satisfactory breakfast. Another approach \DIFdelbegin \DIFdel{we discussed was focusing }\DIFdelend \DIFaddbegin \DIFadd{was focused }\DIFaddend on the cleanliness\DIFdelbegin \DIFdel{and comfortthat Western tourists expect and making the hotels }\DIFdelend \DIFaddbegin \DIFadd{, comfort, and }\DIFaddend tobacco-free \DIFdelbegin \DIFdel{. We favor ``long-haul}\DIFdelend \DIFaddbegin \DIFadd{space expected by Western tourists. ``Long-haul}\DIFaddend '' Western tourists\DIFaddbegin \DIFadd{, }\DIFaddend who spend almost half of their budget on hotels with this strategy\DIFdelbegin \DIFdel{. While }\DIFdelend \DIFaddbegin \DIFadd{, were favored. Although }\DIFaddend Westerners are less in number than Chinese tourists, it could \DIFdelbegin \DIFdel{prove to }\DIFdelend \DIFaddbegin \DIFadd{be proven that they }\DIFaddend have more substantial returns. This is because Chinese customers also favor cleanliness as a satisfaction factor, and both populations could be pleased. This paper provides results and discussion that can be utilized as a guideline for managerial decisions when considering Chinese and Western tourists in Japan. We can observe their stark differences \DIFdelbegin \DIFdel{and }\DIFdelend \DIFaddbegin \DIFadd{as well as }\DIFaddend shared attributes. 

\nobibliography{bibfile-emotion}

\end{document}

